\section{Proofs of statements from earlier work}\label[appendix]{app:earlier}
The following proofs make the supporting interfaces self-contained.
Citations identify the origin of an argument, rather than an assumed
result. The published foundations are those of
\cref{par:five-inputs}; standard field, probability, and arithmetic
facts are used in their usual published forms. In the reproduced
machine constructions, local exponents named \(G\) denote the
general-circuit exponent written \(G^\star\) in \cref{app:transfer};
other variables are local to their stated interfaces.


\subsection{The product-star partition}
\label{earlier:imp:partition}
This is \cite[Lemma~A.1]{Nag72}; the proof is reproduced with notation adapted.

\begin{proof}
Use $rb$ coordinates in $r$ blocks, leaving the others unchanged.
Let $H$ be the $s\times b$ matrix whose columns are the nonzero vectors
of $\F_2^s$. Its rank is $s$, so its kernel $\mathcal C$ has
$2^{b-s}$ elements. Every $x\in\F_2^b$ has exactly one representation
$x=c+e$, with $c\in\mathcal C$ and
$e\in\{0,e_1,\ldots,e_b\}$: the syndrome $Hx$ uniquely specifies $e$.
The same assertion holds with centers in any fixed coset of $\mathcal C$.
Use all $2^s$ cosets in each block. This gives $2^{rs}$ partitions,
and the displayed cardinalities follow by multiplication, including the
unblocked coordinates as center coordinates.

A uniform coset followed by a uniform center is uniform on a block.
Choose all errors independently of both. This proves the averaging
claim. Assign $r/2$ whole blocks to each side. Each contributes exactly
$b-s$ free center bits, so the two blocked parts have equal dimension.
Split the remaining $n-rb$ coordinates as evenly as possible. A fixed
linear basis of each code gives polynomial-time encoding and decoding.
Finally $b+1> (n/r+1)/2$, hence $b+1=\Theta(n/r)$.
\end{proof}

\subsection{Range of the four-corner difference}
\label{earlier:imp:range}
This is \cite[Lemma~A.2]{Nag72}; the proof is reproduced with notation adapted.

\begin{proof}
The underlying real scores at the corners have the form
$a,a+\ell,a+t,a+\ell+t$: the left and right error supports are
disjoint. Their opposite diagonal sums agree. A difference of $2$
would require the two positively signed threshold values in
the defining four-corner difference to be one and the other two to be zero. Their score
sums would then be respectively at least and strictly below twice the
threshold, a contradiction. The case $-2$ is identical with the
diagonals reversed. \end{proof}

\subsection{The influence estimate}
This elementary argument is reproduced from \cite[Lemma~A.3]{Nag72}.

\begin{lemma}[A threshold's total influence]
\label{earlier:star-influence}
If a threshold depends on $f$ input coordinates, then
$\sum_i\Inf_i(T)\le\sqrt f$ under the uniform cube distribution.
\end{lemma}
\begin{proof}
Complement coordinates of negative weight to make the function monotone.
These complements preserve every influence. Use spin variables $X_i$
and the spin-valued output $g=2T-1$. For a monotone function,
conditioning on all coordinates except $i$ shows
$\E[g(X)X_i]=\Inf_i(T)$: the conditional expectation is one exactly
when coordinate $i$ is pivotal and zero otherwise. Therefore
\[
\sum_i\Inf_i(T)=\E\!\left[g(X)\sum_{i\in\supp T}X_i\right]
\le (\E g^2)^{1/2}
     \left(\E\left(\sum_{i\in\supp T}X_i\right)^2\right)^{1/2}
=\sqrt f.
\]
This elementary consequence of unateness is the only influence bound
needed below.
\end{proof}

\subsection{The mixed-dependence population bound}
\label{earlier:imp:population}
This is \cite[Theorem~A.4]{Nag72}; the proof is reproduced with notation adapted.

Write $I_L(T)=\sum_{i\text{ on the left}}\Inf_i(T)$ and let $f_R$
be the right support size. The estimate proved below is
\[
\Pr[\Delta_T\ne0]\le
\frac{2I_L(T)}{b+1}\min\{1,f_R/(b+1)\}
\le\frac{2f}{(b+1)^{3/2}}
\le6\frac{r^{3/2}f}{n^{3/2}}.
\]
\begin{proof}
If $v$ misses the support, the mixed difference is zero. Otherwise,
$\Delta_T\ne0$ requires at least one of the left paths from $c$ to
$c+u$, or from $c+v$ to $c+v+u$, to cross a pivotal edge. Order the
selected left coordinates canonically. After fixing both errors, every
intermediate vertex is a fixed translate of the uniform ambient center,
and so is uniform. The union bound on the two paths is
$2\sum_{i\text{ selected on left}}\Inf_i(T)$.

The errors on the two sides are independent. Each left coordinate is
selected with probability $1/(b+1)$, and a union bound on the right
support gives $\Pr[v\text{ hits support}]\le
\min\{1,f_R/(b+1)\}$. Averaging the path bound proves the first
inequality. Apply \cref{earlier:star-influence} and
$\min\{1,t\}\le\sqrt t$ with $t=f/(b+1)$ for the second.
The last follows from \cref{imp:partition}. Sum over bottom gates;
their support sizes sum to at most the wire count.
\end{proof}

\subsection{A sharp disjoint exact-child decomposition}
\label{earlier:imp:disjoint}
This is \cite[Lemma~C.2]{Nag72}, whose exact-child construction is \cite[Lemma~A.4]{NagCubic}; the proof is reproduced with notation adapted.

\begin{proof} Put \(a_i=|w_i|\), let \(z_i=x_i\) when \(w_i\ge0\) and
\(z_i=1-x_i\) when \(w_i<0\), and set

\[
A_- = \sum_{i:w_i<0}a_i,
\qquad
T'=T+A_-.
\]

The identity

\[
\sum_iw_ix_i=\sum_i a_i z_i-A_-
\]

reduces the gate to \(\mathbf1[\sum_i a_i z_i>T']\). If \(T'<0\), use
the single always-true equation \(0=0\); if \(T'\ge\sum_i a_i\), use the
single always-false equation \(0=1\). Hence assume
\(0\le T'<\sum_i a_i\). Let \(L\) be the least positive integer for
which \(T',a_1,\ldots,a_f\le2^L-1\). Then \(L\le\ell\), and write

\[
a_i=\sum_{b=1}^L2^{b-1}a_{i,b},
\qquad a_{i,b}\in\{0,1\}.
\]

For \((a,b)\in[f]\times[L]\), define the partial score

\[
P_{a,b}(z)
={}
\sum_{i=1}^f z_i
\left(
\sum_{j>b}2^{j-1}a_{i,j}
{}+{}
\mathbf1[i\ge a]\,2^{b-1}a_{i,b}
\right).
\]

Prepend the explicit sentinel \(0\), and then order these scores as

\[
P_{f,L},P_{f-1,L},\ldots,P_{1,L},
P_{f,L-1},\ldots,P_{1,1}.
\]

The sequence is nondecreasing and ends at \(\sum_i a_i z_i\). Let
\(P^-_{a,b}\) denote the entry immediately before \(P_{a,b}\), with the
prepended sentinel used at the first position. Directly from the
definition,

\[
P_{a,b}-P^-_{a,b}=2^{b-1}a_{a,b}z_a.
\]

Consequently, on every accepting input there is a unique first pair
\((a,b)\) whose partial score exceeds \(T'\). At that pair
\(P^-_{a,b}\le T'<P_{a,b}\), so its step is nonzero and hence
\(a_{a,b}=z_a=1\). Write

\[
T'=2^{b-1}Q_b+R_b,
\qquad
0\le R_b<2^{b-1},
\qquad
Q_b\ge0.
\]

Both \(P_{a,b}\) and \(P^-_{a,b}=P_{a,b}-2^{b-1}\) are multiples of
\(2^{b-1}\). Therefore the two inequalities \(P^-_{a,b}\le T'<P_{a,b}\)
hold if and only if \(P_{a,b}=2^{b-1}(Q_b+1)\). Thus \((a,b)\) is the
first crossing if and only if

\[
a_{a,b}=1,
\qquad z_a=1,
\qquad
P_{a,b}(z)=2^{b-1}(Q_b+1).
\]

For a zero bit \(a_{a,b}\) omit the child. For a one bit use the single
exact equation

\[
2P_{a,b}(z)+z_a=2^b(Q_b+1)+1.
\]

Parity in the child equation forces \(z_a=1\), after which the equation is
precisely the crossing conditions. Thus each accepting input satisfies its unique
first-crossing child and every rejecting input satisfies none: the OR is
disjoint.

There are at most \(fL\) nonconstant children. A coefficient on the left
of the child equation is at most \(2(2^L-1)+1=2^{L+1}-1\). Since
\(Q_b\le2^{L-b+1}-1\), its right-hand side is at most \(2^{L+1}+1\).
Substitution of \(z_i=1-x_i\) negates the corresponding coefficient and
subtracts it from the target, giving coefficients of magnitude at most \((f+2)2^{\ell+1}\), hence the \(O(\ell+\log(f+2))\) bits of \cref{imp:disjoint}. Finally, in the nonconstant case
\(T'<\sum_{i=1}^f a_i\le fB_0'\), so \(L\le\ell\); all loops and
arithmetic are polynomial in the displayed output size. \end{proof}

\subsection{Separated pullback through the center coordinates}
\label{earlier:imp:pullback}
This is \cite[Lemma~C.3]{Nag72}; the proof is reproduced with notation adapted.

\begin{proof}
Write an original child as $c_j+\sum_i d_{j,i}x_i=0$.  Put its constant and the terms whose coordinate depends on $z_{\rm L}$ into $a_j(z_{\rm L})$, and the remaining terms into $b_j(z_{\rm R})$.  Evaluating at most $n$ coordinate
functions and adding polynomial-bit integers takes polynomial time and
increases bit length by at most $O(\log n)$.  The identity and disjointness
hold for every original Boolean input, so substitution preserves both.

\end{proof}
In the product-star parametrization, keep each entire code block on one
column side.  Each coordinate of its center, including the subsequent
exclusive-or with a fixed error row, then depends on that column side
alone.  \Cref{imp:pullback} applies even though a center bit may be a
parity of many message bits.  It does \emph{not} assert that a real
threshold remains a real threshold under a linear map over
$\mathbb F_2$.  



\subsection{Stacking equations}
\label{earlier:imp:stack}
This is \cite[Lemma~C.6]{Nag72}; the proof is reproduced with notation adapted.

\begin{proof}
The forward implication is immediate.  Otherwise choose the largest $j$
with $E_j(z)\ne0$.  Its contribution has magnitude at least $H^{j-1}$,
whereas all lower contributions have total magnitude at most
$V(H^{j-1}-1)/(H-1)<H^{j-1}$.  Cancellation is impossible.
Sum the left half-scores and the right half-scores separately with the
same powers of $H$.  Since $H=2V+2$, the largest power has
$O(t\log(V+2))$ bits, proving the bounds.  For the empty conjunction,
use the equation $0=0$.
\end{proof}

\paragraph{The published multiplication input.}
\begin{imported}[Williams]\label{earlier:rectangular-input}
There is one deterministic multitape Turing machine that, for all
sufficiently large \(N\), multiplies over the integers any two
\(0\)-\(1\) matrices of dimensions \(N\times\ceil{N^{0.1}}\) and
\(\ceil{N^{0.1}}\times N\) in \(N^2\poly(\log N)\) steps.
\end{imported}
This is \cite[Corollary~4.4 and Appendix~C]{WilACC14}.

\subsection{Signed rectangular multiplication}
This is \cite[Lemma~C.4]{Nag72}; the proof is reproduced with notation adapted.

\begin{lemma}[Signed-by-Boolean multiplication]
\label{earlier:alg-rectangular}
Fix $C$.  Suppose $d\le\lceil U^{0.1}\rceil$,
$A\in\mathbb Z^{U\times d}$, and $B\in\{0,1\}^{d\times U}$.
If the entries of $A$ have at most $(\log_2(U+2))^C$ bits, then the exact
product $AB$ takes $U^2\poly(\log U)$ bit operations.
\end{lemma}

\begin{proof}
Let $b$ bound the magnitude bit length, and write
$A=\sum_{t=0}^{b-1}2^t(A_t^+-A_t^-)$, using separate Boolean magnitude
planes for positive and negative entries.  Apply \cref{earlier:rectangular-input} to
the $2b$ products $A_t^\pm B$, padding the middle dimension with zeros.
Their entries lie in $[0,d]$, so signed recombination uses
$O(b+\log(d+1))$ bits per answer and the asserted total time.  This is
the binary-slicing reduction of \cite[Lemma~B.2]{NagCubic}.
\end{proof}

\subsection{The separated exact printer}
\label{earlier:imp:printer}
This is \cite[Theorem~C.5]{Nag72}; the proof is reproduced with notation adapted.

\begin{proof}
This is the weighted-dominance proof of
\cite[Theorem~B.3]{NagCubic}, with arbitrary evaluable half-scores in
place of linear half-scores.  We give the reduction to verify that this is
the only property of linearity it uses.  Replace an equality by the two
successor cuts
\[
 \ind[E_i=0]=\ind[E_i\ge0]-\ind[E_i\ge1].
\]
For even $n_{\mathrm{out}}$, put $U=2^{n_{\mathrm{out}}/2}$.  Enumerate the half-scores of the resulting
$m\le2s$ cuts to express the required table as
\[
 P[a,b]=\sum_{j=1}^m\omega_j\ind[A[a,j]\le B[j,b]].
\]
This takes $Um\poly(n)$ time.  For all sufficiently large $n_{\mathrm{out}}$,
$m\le2U^{0.02}\le U^{0.049}$.  Pad the gate dimension to
$\lceil U^{0.049}\rceil$ using zero coefficients.

For each $j$, sort the $2U$ labeled entries in its $A$ column and $B$
row, putting every $A$ entry before every $B$ entry at an equal numerical
value.  Thus a weak comparison is true exactly when the $A$ rank is
earlier.  Partition the ordered list into consecutive buckets of at most
$\lceil U^{0.951}\rceil$ entries.  For each pair $(a,j)$, inspect the
$B$ entries in the same bucket, and add $\omega_j$ precisely when their
ranks are ordered correctly.  There are $O(U^2)$ such inspections.

For the remaining comparisons, form
\[
 A'[a,(j,h)]=\omega_j\ind[A[a,j]\text{ is in bucket }h],
 \qquad
 B'[(j,h),b]=\ind[B[j,b]\text{ is in a later bucket than }h].
\]
The common dimension is $O(U^{0.098})$, hence at most
$\lceil U^{0.1}\rceil$ eventually.  \cref{earlier:alg-rectangular}
computes $A'B'$ in $U^2\poly(n)$ bit operations.  Its entries are exactly
the cross-bucket contributions.  Add the already computed same-bucket
contributions.  All sorting, table construction, and additions fit the
same bound.  Coefficient and output lengths remain polynomial because
each output is a sum of at most $2s$ polynomial-bit integers.

For odd $n_{\mathrm{out}}$, add one ignored bit to the smaller half, run the even case,
and retain its zero slice.  This multiplies time by at most two.
The finitely many sizes below the onset of \cref{earlier:rectangular-input} can be treated directly.
At no point was a half-score required to be linear in its argument.
\end{proof}

\subsection{A short seed and exact conditional expectations}
\label{earlier:imp:seed}
This is \cite[Lemma~F.2]{Nag72}; the proof is reproduced with notation adapted.

\begin{proof}
Use the standard polynomial-evaluation construction of limited
independence, in the universal-hashing tradition of
Carter--Wegman~\cite{CarterWegman}; its needed properties are verified
here. Take $\ell\ge\max\{k,\lceil\log_2(Q+1)\rceil\}$ and a field
$\F_{2^\ell}$. Label the $Q$ samples by distinct field elements.
Evaluate a uniformly random polynomial of degree $2h-1$ at those labels
and project each value to $k$ binary coordinates. Evaluations at any
$2h$ distinct labels are independent uniform field elements, so the
projected positions are $2h$-wise independent and uniform on $[N]^2$.
The seed consists of $2h\ell$ binary coefficient bits. A defining
irreducible polynomial and the arithmetic can be found and implemented
in $2^{O(\ell)}$ elementary preprocessing time, which is allowed below;
no randomness in field construction is needed.

For a fixed prefix, every assertion that a sample lies at a specified
position is a system of binary linear equations in the remaining seed
bits. A product of such indicators is their conjunction. Gaussian
elimination gives conditional probability zero if it is inconsistent,
and $2^{-a}$ if its remaining rank is $a$. The denominator divides
$2^S$. The trace expansion is finite, so these probabilities compute
all conditional coefficients exactly. Repeated sample labels simply
repeat or augment the same system.

The averaging identity
is the ordinary partition of the remaining uniform seed space into its
two equally sized children. 
\end{proof}

\subsection{Recovering distributions from modular counts}\label{earlier:g:sec:windows}\label{earlier:gate-proof-group}\label{earlier:gate-window}

This reproduces \cite[Section~4]{NagCG}, with notation adapted. The counter of \cref{sec:printing} will supply exact residue counts of
integer scores. Residues must determine a \emph{distribution}, not just
a single integer. For example, knowing how often a score is even does not
identify which even scores occurred. Our reconstruction lemmas apply
to arbitrary rational-valued functions; positivity is not required.
This permits a clean uniqueness proof by subtracting two candidate
distributions.

Let \(\cP\) be a nonempty finite set of primes and put
\[
 D'=1+\sum_{p\in\cP}(p-1).
\]
For a finitely supported function \(F:\Z^d\to\Q\), define its
modular push-forward by
\[
 F_p(\rho)=\sum_{u\equiv\rho\pmod p}F(u),
 \qquad\rho\in(\Z/p\Z)^d.
\]
Thus \(F_p\) records the total mass of each residue class.

\subsection{One score: a roots-of-unity argument}

\begin{lemma}\label{earlier:g:lem:one-window}
If \(F:\Z\to\Q\) is supported on at most \(D'\) consecutive
integers and \(F_p=0\) for every \(p\in\cP\), then \(F=0\).
\end{lemma}
\begin{proof}
Shift the support to \(\{0,\ldots,D'-1\}\) and form
\(f(t)=\sum_uF(u)t^u\). For every \(p\)-th root of unity \(\zeta\),
\[
 f(\zeta)=\sum_{\rho=0}^{p-1}F_p(\rho)\zeta^\rho=0.
\]
Distinct primes' root sets intersect only at one. Their union therefore
contains \(D'\) distinct roots. A polynomial of degree below \(D'\)
with this many roots is zero.
\end{proof}

\subsection{Two scores: short projections and binomial factors}

Let \(\varphi\) be Euler's totient function and set
\(S(q)=\sum_{m=1}^q m\varphi(m)\), with \(S(0)=0\).
Hajdu and Tijdeman's discrete-tomography method~\cite[Theorem~1 and
its preliminary divisibility lemma]{HT01} turns zero line sums into
binomial divisibility. We prove this over \(\Q\) and count the usable
directions.

\begin{lemma}\label{earlier:g:lem:two-window}
Let \(L\ge2\) and
\(q=\floor{(D'-1)/(L-1)}\). If \(S(q)\ge L\), then a
rational-valued function supported on a given box with both sides at
most \(L\) is uniquely determined, among functions supported on that
box, by its push-forwards \(F_p\), \(p\in\cP\).
\end{lemma}

\emph{Idea of proof.} Take the difference of two functions with the
same data. Each of its projections along a short primitive direction is
a one-score function with zero data, hence zero by the one-score lemma.
Each vanishing projection makes an irreducible binomial divide the
generating polynomial. These binomials have more total degree in one
variable than the polynomial has, so the difference is zero.

\begin{proof}
Subtract two functions with the same data and translate their common
box to \(\{0,\ldots,L-1\}^2\). It suffices to show that a function
with zero data is zero. Write its generating polynomial as
\[
 f(u,v)=\sum_{a,b=0}^{L-1}F(a,b)u^av^b.
\]

\emph{First, every short primitive projection is zero.}
A vector \((\alpha,\beta)\ne(0,0)\) is primitive if
\(\gcd(\abs\alpha,\abs\beta)=1\). If
\(\abs\alpha+\abs\beta\le q\), then
\[
 G(c)=\sum_{\alpha a+\beta b=c}F(a,b)
\]
is supported on an interval containing at most
\((L-1)q+1\le D'\) integers. Its residue data are zero, since
projection commutes with reduction modulo \(p\). By
\cref{earlier:g:lem:one-window}, \(G=0\). Equivalently,
\(f(t^\alpha,t^\beta)=0\) in \(\Q[t,t^{-1}]\).

\emph{Second, each vanishing projection forces a distinct prime factor.}
Choose one vector from each pair \(\pm(\alpha,\beta)\), normalized so
that \(\alpha\ge0\), and \(\beta=1\) if \(\alpha=0\).
The kernel of the substitution \(u\mapsto t^\alpha,v\mapsto t^\beta\)
on \(\Q[u,v]\) is generated by the binomial
\[
 B_{\alpha,\beta}=
 \begin{cases}
 u^\beta-v^\alpha,&\beta\ge0,\\
 u^{\abs\beta}v^\alpha-1,&\beta<0.
 \end{cases}
\]
In the Laurent polynomial ring,
choose integers \(i,j\) with \(\alpha i+\beta j=1\).
The monomials \(\omega_1=u^iv^j\) and \(\omega_2=u^\beta v^{-\alpha}\)
form a coordinate system, because their exponent matrix has
determinant \(-1\). Substitution sends \(\omega_1\) to \(t\) and
\(\omega_2\) to one, so its kernel there is \((\omega_2-1)\). Up to a Laurent monomial unit, this is
the displayed binomial. That binomial has no factor \(u\) or \(v\).
It is irreducible in \(\Q[u,v]\): a factorization gives a Laurent
unit factor, since \(\omega_2-1\) is irreducible in the Laurent ring,
hence a monomial, which must be constant.
Clearing Laurent denominators and using unique factorization gives
the asserted polynomial kernel.

Nonparallel primitive directions give nonassociate binomials:
substitution in the other direction gives two powers of \(t\)
whose exponents differ by \(\pm(\alpha\beta'-\beta\alpha')\ne0\).
By unique factorization, the product of these binomials divides \(f\).

\emph{Finally, there are too many factors for a nonzero \(f\).}
The sum of their degrees in \(u\) is
\(\sum\abs\beta\), over primitive directions modulo sign of norm
at most \(q\). At norm one, the horizontal and vertical directions
contribute one in total. At norm \(m\ge2\), the directions contribute
\[
 2\sum_{\substack{1\le a<m\\\gcd(a,m)=1}}(m-a)
 =m\varphi(m),
\]
using the symmetry \(a\leftrightarrow m-a\). Thus the product has
\(u\)-degree \(S(q)\ge L\), whereas \(\deg_uf\le L-1\).
This is impossible unless \(f=0\).
\end{proof}

For a concrete picture, the projection \((\alpha,\beta)=(1,1)\)
adds masses along diagonals \(a+b=c\). If all those sums vanish,
\(u-v\) divides \(f\). The projection \((1,-1)\) instead forces
\(uv-1\). The lemma uses many such independent factors, not just
the coordinate projections.

\begin{lemma}\label{earlier:g:lem:window-size}
For every integer \(q\ge1\),
\[
 S(q)\ge\frac{2q^3}{\pi^2}-q^2(4/3+\ln q).
\]
Call \(L\ge2\) admissible if \(S(\floor{(D'-1)/(L-1)})\ge L\), as in
\cref{earlier:g:lem:two-window}. Let \(L_2\) be the largest admissible \(L\),
or one if none exists. Then
\(L_2=\Theta(D'^{3/4})\) as \(D'\to\infty\).
In particular \(L_2\ge\floor{(3/5)D'^{3/4}}\) for all sufficiently
large \(D'\).
\end{lemma}
The proof is a standard totient-summation calculation, given in
\cref{earlier:g:app:window-size}.

Uniqueness also gives an algorithm: introduce one rational unknown
per box point and one equation per residue class. The matrix entries
are zero or one. Exact Gaussian elimination finds the unique
distribution whenever the data come from a function on the box.
An algorithm may also detect inconsistency and report uncertainty.
A box with one point is determined by its total mass, so the
singleton convention in the definition of \(L_2\) is also sound.

\subsection{How many useful primes are available?}

For an integer \(P\ge3\), let \(\cP\) consist of all primes in
\((P/2,P]\). Only the order of magnitude of this set is needed.

\begin{lemma}\label{earlier:g:lem:primes}
For this prime set, as \(P\to\infty\),
\[
 D'=\Theta(P^2/\ln P),\qquad
 L_2=\Theta\!\left(P^{3/2}/(\ln P)^{3/4}\right).
\]
\end{lemma}
\begin{proof}
By Chebyshev-type bounds (for instance Erd\H{o}s' proof of
Bertrand's postulate~\cite{Erdos32}), the interval \((P/2,P]\)
contains \(\Theta(P/\ln P)\) primes.
Each such prime contributes
more than \(P/2-1\) and at most \(P\) to \(D'\), so
\(D'=\Theta(P^2/\ln P)\). Apply \cref{earlier:g:lem:window-size}.
\end{proof}

\begin{remark}[The limitation of full-box residue reconstruction]\label{earlier:g:rem:rank}
For distributions on \(\{0,\ldots,L-1\}^d\), the residue map has
rank at most \(1+\sum_{p\in\cP}(p^d-1)\): all primes report the
same total mass. Injectivity therefore requires
\(L^d\le1+\sum_p(p^d-1)\).
For two scores this yields
\(L=O(P^{3/2}/\sqrt{\ln P})\). Our box bound has the same power
of \(P\), with an additional factor \((\ln P)^{1/4}\) loss.
This is a restriction on recovering \emph{arbitrary} box-supported
distributions from these data, not an impossibility theorem for all
circuit-analysis methods or for more structured distributions.
\end{remark}



\subsection{Totient sums and the two-score window}\label{earlier:g:app:window-size}\label{earlier:gate-window-size}
\begin{proof}[Proof of \cref{earlier:g:lem:window-size}]
This reproduces \cite[Appendix~A]{NagCG}, with notation adapted. These are standard totient-summation calculations; see
Apostol~\cite[Chapters~2--3 and~11]{Apostol76} for the M\"obius
identities, arithmetic averages, and Euler products.
Write \(\mu\) for the M\"obius function. Since
\(m\varphi(m)=\sum_{d\mid m}\mu(d)m^2/d\),
\[
 S(q)=\sum_{d\le q}\mu(d)d\sum_{e\le q/d}e^2.
\]
The sum-of-squares formula gives
\(\abs{\sum_{e\le y}e^2-y^3/3}\le y^2\) for \(y\ge1\).
Therefore
\[
 S(q)\ge\frac{q^3}{3}\sum_{d\le q}\frac{\mu(d)}{d^2}
          -q^2\sum_{d\le q}\frac1d.
\]
Use \(\sum_{d\ge1}\mu(d)/d^2=6/\pi^2\), a tail of magnitude at most
\(1/q\), and \(\sum_{d\le q}1/d\le1+\ln q\).
For \(L=\floor{cD'^{3/4}}\), the associated \(q\) is
\((1/c+o(1))D'^{1/4}\). Thus \(S(q)\ge L\) eventually whenever
\(c^4<2/\pi^2\); this includes \(c=3/5\).
Conversely \(S(q)\le\sum_{m\le q}m^2=O(q^3)\), and
\(q\le(D'-1)/(L-1)\), so every admissible \(L\) is
\(O(D'^{3/4})\).
\end{proof}




\subsection{Nested subsampling produces short integer digits}\label{earlier:g:sec:chains}\label{earlier:gate-chain}

This reproduces \cite[Section~5]{NagCG}, with notation adapted. We now turn a potentially large live sum into short forms that can
be identified from one residue. Sampling at a single very small rate
would make the first sampled sum short, but leave a poor estimate of
the original sum. The chain increases the rate while keeping each
new correction short. Integer rate ratios preserve exact arithmetic.

\begin{lemma}[Choice of rates]\label{earlier:g:lem:chain}
Let \(T>0\), let \(h\) be an integer with \(h\ge64T\), and let
\(\ell_0\ge0\). Define
\[
 \bar\ell=\max\{\ell_0,h^2/(256T)\},\quad
 \pi^*=\frac{h^2}{8T\bar\ell},\quad
 s(\pi)=\sqrt{2T\pi\bar\ell}+\frac{2T}{3},
\]
\[
 M_1=\ceil{4\bar\ell/h},\qquad \pi_1=1/M_1.
\]
While \(x_j=\pi_j/\pi^*<2^{-8}\), set
\[
 a_j=\floor{1/(4\sqrt{x_j})},\qquad \pi_{j+1}=a_j\pi_j.
\]
Stop at the first index \(J\) with \(x_J\ge2^{-8}\). Then
\begin{enumerate}[label=(\roman*)]
\item \(a_j\ge4\), all rates lie in \((0,1]\), and
\(\pi_j=\prod_{i<j}a_i/M_1\).
\item \(\pi_1\bar\ell+s(\pi_1)\le h/2\), and
\(s(\pi_{j+1})+a_js(\pi_j)\le h/2\) for \(j<J\).
\item \(J\le\ceil{\lgx\lgx(h/T)}+1\).
\item \(s(\pi_J)/\pi_J\le86T\bar\ell/h\).
\end{enumerate}
\end{lemma}
\begin{proof}
Since \(h\ge64T\), the second entry of the maximum gives
\(\bar\ell\ge h^2/(256T)\ge h/4\).
Thus \(4\bar\ell/h\le M_1\le8\bar\ell/h\), and
\[
 T/h\le x_1\le2T/h\le1/32.
\]
Write \(\tau(\pi)=\sqrt{2T\pi\bar\ell}
=(h/2)\sqrt{\pi/\pi^*}\).
For the first digit,
\(\pi_1\bar\ell\le h/4\),
\(\tau(\pi_1)\le\sqrt{Th/2}\le h/8\), and \(2T/3\le h/96\).

During a step, \(a_j\ge4\),
\(a_j\sqrt{x_j}\le1/4\), and
\(a_jx_j\le\sqrt{x_j}/4\), so \(\sqrt{a_jx_j}\le x_j^{1/4}/2<1/8\). Consequently
\[
 \tau(\pi_{j+1})+a_j\tau(\pi_j)
 =\frac h2\bigl(\sqrt{a_jx_j}+a_j\sqrt{x_j}\bigr)
 \le\frac{3h}{16}.
\]
The \(x_j\) increase, so
\(a_j\le\sqrt{h/T}/4\), and
\[
 (1+a_j)\frac{2T}{3}
 \le h/96+\sqrt{hT}/6\le h/32.
\]
These bounds prove (ii).
Every newly produced \(x_{j+1}\) is below \(1/64\), while
\(x_1\le1/32\). Also \(\pi^*\le32\), so all rates are at most one.

Since \(1/(4\sqrt{x_j})>4\) during a step, flooring loses at most
a factor two, and
\(x_{j+1}\ge\sqrt{x_j}/8\).
For \(y_j=\lgx(1/x_j)\), this means
\(y_{j+1}-6\le(y_j-6)/2\).
If the chain has not stopped earlier, iteration gives
\(y_j\le6+2^{-(j-1)}y_1\), with
\(y_1\le\lgx(h/T)\). A step at index \(j\) requires \(y_j>8\), hence
\(2^{j-1}<y_1/2\) and \(j<\lgx y_1\). Applied to \(j=J-1\) when
\(J\ge2\), this gives \(J<\lgx\lgx(h/T)+1\), which proves (iii).
At the stopping index, \(x_J\ge2^{-8}\), so
\(\pi_J\ge h^2/(2048T\bar\ell)\); with \(T/h\le1/64\) this gives
\[
 \frac{s(\pi_J)}{\pi_J}
 \le\frac{64T\bar\ell}{h}
       +\frac{1366T^2\bar\ell}{h^2}
 \le\frac{86T\bar\ell}{h}.
\]
This proves (iv) and completes the argument.
\end{proof}

Expand a top sum \(X\) into signed unit copies \(\sigma_tg_{i(t)}\).
Choose independent uniform \(u_t\in[0,1]\) and define
\[
 K_j=\sum_t\ind[u_t\le\pi_j]\sigma_tg_{i(t)},\qquad
 E_0=K_1,\quad E_j=K_{j+1}-a_jK_j\ (1\le j<J).
\]
Let
\[
 A=\prod_{i=1}^{J-1}a_i,\qquad A_0=A,\qquad
 A_j=\prod_{i=j+1}^{J-1}a_i\quad(1\le j<J).
\]
The identities
\begin{equation}\label{earlier:g:eq:telescope}
 K_J=\sum_{j=0}^{J-1}A_jE_j,\qquad \pi_J=A/M_1
\end{equation}
follow by repeatedly substituting \(K_{j+1}=a_jK_j+E_j\).
The first identity holds for \emph{any} proposed coefficient vectors
\(K_1,\ldots,K_J\), whether or not they arose from nested sampling.
\Cref{prop:sound} uses this.



\subsection{The Beigel--Tarui lifting polynomial}
\label{earlier:gate-amp}
This reproduces \cite[Lemma~6.6]{NagCG}, with notation adapted.

\begin{lemma}\label{earlier:g:lem:amp}
For \(e\ge1\), define
\[
 \Amp_e(t)=1-(1-t)^e\sum_{j=0}^{e-1}\binom{e-1+j}{j}t^j.
\]
This integer polynomial has degree \(2e-1\), with
\(t^e\mid\Amp_e(t)\) and \((1-t)^e\mid1-\Amp_e(t)\).
Consequently, modulo \(p^e\), it maps an integer congruent to zero
modulo \(p\) to zero, and one congruent to one to one.
\end{lemma}
This is the modulus-amplifying polynomial of Beigel and
Tarui~\cite[Section~2.3]{BT94}.
\begin{proof}
The power-series identity
\((1-t)^{-e}=\sum_{j\ge0}\binom{e-1+j}{j}t^j\) shows that truncating
at degree \(e-1\) leaves an error divisible by \(t^e\).
Multiplication by \((1-t)^e\) proves the first divisibility assertion.
The second is immediate from the definition. Both imply the stated
congruences, including for negative integer arguments.
\end{proof}



\subsection{Pointwise proximity and the five-weight gadget}
\label{earlier:gate-pcpp-gadget}
This reproduces \cite[Sections~11.1--11.3]{NagCG}, with notation adapted.

\begin{imported}[Pointwise three-query PCPP]\label{earlier:g:imp:three-pcpp}
There are a rational constant $0<\gamma<1$, a polynomial-time
computable linear encoding $\operatorname{Enc}$, and a polynomial-time
transformation taking a Boolean circuit $D$ on $q$ inputs to a
polynomial-size list of three-literal clauses on variables $Y\cup Z$,
such that the following hold for every $x\in\{0,1\}^q$ with
$Y=\operatorname{Enc}(x)$:
\begin{itemize}
\item If $D(x)=1$, an auxiliary assignment computable in polynomial
time satisfies every clause.
\item If $D(x)=0$, every auxiliary assignment leaves at least a
$\gamma$ fraction of the clauses false.
\end{itemize}
Each systematic coordinate is a parity of input coordinates. The
number of clauses can be taken to be a power of two.
\end{imported}

This is the pointwise form of Chen--Williams' perfectly complete
three-query PCPP~\cite[Lemma~24 and Appendix~A]{CW19}, using their
linear code and decoder~\cite[Lemma~27]{CW19}, which \cref{earlier:release-interfaces} derives from~\cite[Lemma~26]{CW19}. As in the proof of
\cite[Theorem~5]{CW19}, here is the reduction to the stated form. Take a fixed proximity radius smaller than the
code's decoding radius, totalize the decoder by returning an arbitrary
fixed string on failure, and apply the PCPP to
$\widetilde D(y)=D(\operatorname{Dec}(y))$. If $D(x)=0$, every word sufficiently
close to $\operatorname{Enc}(x)$ decodes to $x$, and is therefore
rejected by $\widetilde D$. Thus $\operatorname{Enc}(x)$ is farther than the
proximity radius from every accepting input of $\widetilde D$, and proximity
soundness applies. If $D(x)=1$, perfect completeness and the efficient
honest-proof procedure of~\cite[Lemma~24]{CW19} apply. The underlying proximity construction
is due to Ben-Sasson, Goldreich, Harsha, Sudan, and
Vadhan~\cite{BGHSV06}. In the constant-size clause conversion in
the proof of \cite[Lemma~24]{CW19}, pad to a power-of-two number of
clauses with tautologies before selecting a clause uniformly. This
preserves perfect completeness and changes only the fixed rejection
constant. Finally choose a smaller rational $\gamma$ if necessary.

For one three-literal clause, let $x_1,x_2,x_3$ be its literal values.
The classical Garey--Johnson--Stockmeyer gadget~\cite{GJS76}, in
the form recorded by Chen--Williams~\cite[Lemma~47]{CW19}, introduces
a Boolean helper $y$ and uses the ten constraints
\[
 x_1,\ x_2,\ x_3,\quad
 \neg x_1\vee\neg x_2,\
 \neg x_2\vee\neg x_3,\
 \neg x_3\vee\neg x_1,\quad
 y,\quad x_1\vee\neg y,\ x_2\vee\neg y,\ x_3\vee\neg y.
\]
If $k=x_1+x_2+x_3$, their total satisfaction score is
$\operatorname{sat}(x,y)=6+k-\binom{k}{2}+y(k-2)$. We use the penalty
\begin{align}
 Q(x_1,x_2,x_3,y)
 &=7-\operatorname{sat}(x,y)\notag\\
 &=1-\sum_i x_i+\sum_{i<j}x_ix_j+2y-y\sum_i x_i.
 \label{earlier:g:eq:gadget-penalty}
\end{align}
Its values on Boolean arguments are
\[
\begin{array}{c|rrrr}
 k&0&1&2&3\\ \hline
 Q(k,0)&1&0&0&1\\
 Q(k,1)&3&1&0&0
\end{array}
\]
Thus $Q\ge0$, and a false original clause forces $Q\ge1$ for
either helper. For a true original clause, the choice
$y^*=x_1x_2x_3$ gives penalty zero.

\begin{lemma}[Weighted stability of the gadget]
\label{earlier:g:lem:gadget-stability}
For every true honest triple $x\ne000$, with $y^*=x_1x_2x_3$, and
every Boolean replacement $x',y'$, we have
\[
 Q(x',y')\le\sum_{i=1}^3|x_i-x_i'|+2|y^*-y'|.
\]
\end{lemma}
\begin{proof}
If $y'=y^*=0$, the honest triple has weight one or two. The only
nonzero penalties with helper zero occur at weights zero and three,
each requiring an input change. If $y'=y^*=1$, the honest triple
is $111$, and the table gives $Q(k',1)\le3-k'$, its Hamming distance
from that triple. If the helper changes, its weight already pays two.
The only penalty exceeding two is $Q(0,1)=3$; this also requires at
least one input change from a true honest triple.
\end{proof}

The aggregate seven-versus-six score is classical~\cite{GJS76}, and
Chen--Williams use it in this aggregate form to prove
\cite[Lemma~25]{CW19}; what is added here is its weighted stability
under proof-table errors. This estimate uses cancellation within the
ten-clause score. Applying
a separate union bound to every clause would discard that cancellation.
The helper has weight two, and each original coordinate has weight one.

\subsection{A certified upper estimate using pair means}

Writing the two-literal clauses in the XOR basis, as Chen--Williams
do for two-query predicates~\cite[proof of Theorem~5]{CW19}, the
identity $ab=(a+b-(a\oplus b))/2$ gives
\begin{equation}\label{earlier:g:eq:gadget-pair-identity}
 Q=1+\frac12\left(
 y-\sum_i x_i-\sum_{i<j}(x_i\oplus x_j)
                 +\sum_i(x_i\oplus y)\right).
\end{equation}
There are ten probability terms, all with coefficient of magnitude
$1/2$. A literal is a guessed Boolean proof coordinate or its
complement, or a prescribed systematic parity or its complement. Hence every nonexact
term is a single-circuit, pair-XOR, or circuit-XOR-raw-parity call.
Pure parity terms, constants, and repeated-coordinate identities are
evaluated exactly. No closure of the target class under XOR is used.

For each nonexact term let $(\widehat p,\widehat\kappa)$ be the
counter's estimate and certified radius, and let $\sigma\in\{-1,1\}$
be its sign in \cref{earlier:g:eq:gadget-pair-identity}. Substitute
$\sigma\widehat p+\widehat\kappa$ for the signed probability and
average the resulting expressions over all gadgets. Call the result
$U$, and write $\mathbb E Q$ for the exact average of the penalty over
gadgets and over $x\in\{0,1\}^q$, evaluated on the guessed coordinates.
Soundness of the intervals gives
\begin{equation}\label{earlier:g:eq:gadget-upper}
 U\ge\mathbb E Q\quad\text{for every witness}.
\end{equation}
For a useful witness (\cref{prop:complete}) with all radii at most
$\omega$, it also gives
\begin{equation}\label{earlier:g:eq:gadget-upper-complete}
 U\le\mathbb E Q+10\omega.
\end{equation}
Indeed a signed probability endpoint loses at most $2\omega$, and
its coefficient has magnitude $1/2$. Arbitrary witnesses may return
broader intervals; only the existence of a useful witness asserts the
small-radius bound.



\subsection{The succinct source construction}\label{earlier:source-construction}
This is \cite[Section~11.4]{NagCG}; the construction is reproduced
with notation adapted. The general-circuit exponent is \({G^\star}\)
as in \cref{sec:ae-constants}; \(h_{\rm sim}=6\) by
\cref{lem:ae-wires}. The derived short-PCP interface is
\cref{imp:transfer-short-pcp}.

For a general circuit $B$ with
$a\ge\nu$ inputs, let $\mathrm{tt}_N(B)\in\{0,1\}^N$ be the word whose
$i$th bit is $B(\mathrm{bin}_\nu(i)0^{a-\nu})$, where
$\mathrm{bin}_\nu(i)$ is the $\nu$-bit binary expansion of $i<N$. Fix a
self-delimiting binary description of general circuits that lists the
gates, so that a description of $n$ bits has at most $n$ gates, and
whose length is polynomial in the numbers of gates and inputs. Fix an integer
$G'\ge2$, depending only on ${G^\star}$ and on this description format, such
that for all large $\nu$ every description of a circuit with at most
$q^{G^\star}$ gates on $q$ inputs has fewer than $\nu^{G'}$ bits. Such a $G'$
exists: these descriptions have length at most a fixed polynomial in
$q=O(\nu)$ whose degree depends only on ${G^\star}$ and the format, and any $G'$
above that degree works. Put
\[
 L_\nu=\nu^{G'},
\]
and, when the description of $B$ has at most $L_\nu$ bits, let
$\langle B\rangle_\nu\in\{0,1\}^{L_\nu}$ be that description followed by
zeros.

\Cref{earlier:g:lem:succinct-verifier} in \cref{earlier:g:app:multitape} gives one
deterministic multitape verifier $\mathcal V_{{G^\star}}(u,y)$, a nondecreasing
time bound $T_{{G^\star}}(n)\ge n$ within which $\mathcal V_{{G^\star}}$ halts on every input
of length $n$ and every witness, and a constant $c_{{G^\star}}$. They satisfy
$T_{{G^\star}}(L_\nu)=c_{{G^\star}}\Tclk(N)$, and, for every such $B$, some witness of length
$T_{{G^\star}}(L_\nu)$ makes $\mathcal V_{{G^\star}}$ accept $\langle B\rangle_\nu$ exactly
when $\mathrm{tt}_N(B)\in L_{\Tclk}$. The only size bound that
$\mathcal V_{{G^\star}}$ places on $B$ is the length $L_\nu$, which depends on
$N$ and the fixed $G'$ alone; the weak machine enforces the bound $q^{G^\star}$
separately. Only three properties of $L_\nu$ are used:
\begin{itemize}
\item it is a fixed polynomial in $\nu$;
\item it has degree at least two;
\item it eventually exceeds the description length of every circuit
with at most $q^{G^\star}$ gates on $q$ inputs.
\end{itemize}
$L_\nu$ is fixed before $\mathcal V_{{G^\star}}$, and these
properties hold whatever constants are chosen after it.

We apply \cref{imp:transfer-short-pcp} once, to $\mathcal V_{{G^\star}}$ with the time bound
$T_{{G^\star}}$, which is computable from the input length in polynomial time, as
that input requires (\cref{earlier:g:lem:succinct-verifier}(i)). On an input $\langle B\rangle_\nu$ it constructs, in time
$\poly(L_\nu,\log T_{{G^\star}}(L_\nu))=\nu^{O_{{G^\star},C_{\rm clk}}(1)}$, an oracle circuit
$C^O_{\langle B\rangle_\nu}$ of size $\nu^{O_{{G^\star},C_{\rm clk}}(1)}$ with
$\ell_{\rm out}=\lgx\Tclk(N)+O(\log\log\Tclk(N))$ inputs. Pad it to the
exact native width
\[
 q(N)=\lceil\log_2\Tclk(N)\rceil+
 A_{\rm pad}\lceil\log_2(\log_2\Tclk(N)+2)\rceil+A_{\rm pad},
\]
by adding inputs that it ignores, for a sufficiently large fixed integer
$A_{\rm pad}$, chosen after $C_{\rm clk}$ as at the end of
\cref{earlier:g:app:multitape} so that $\ell_{\rm out}\le q(N)$. Then
$q=\Theta(\log N)$ and $2^q/\Tclk(N)=q^{O(1)}$.


\paragraph{The outer verifier.}
The source-circuit test is
\begin{equation}\label{earlier:g:eq:source-check}
 B_{\rm src}\bigl(\mathrm{bin}_\nu(i)0^{q-\nu}\bigr)=\zeta_i
 \qquad\text{for every }i<N,
\end{equation}
that is, $\mathrm{tt}_N(B_{\rm src})=\zeta$, and rejects otherwise. It
constructs $C^O_{\langle B_{\rm src}\rangle_\nu}$, whose number
$\ell_{\rm out}$ of inputs may depend on $B_{\rm src}$, and guesses a
general circuit $B_{\rm pf}$ with $q$ inputs and at most $q^{G^\star}$ gates,
the outer proof oracle, again rejecting malformed descriptions. It
answers an oracle query $a\in\{0,1\}^{\ell_{\rm out}}$ by
\[
 O_{B_{\rm pf}}(a)=B_{\rm pf}\bigl(a0^{q-\ell_{\rm out}}\bigr),
\]
the same zero padding as in \eqref{earlier:g:eq:source-check}. Write
$B=(B_{\rm src},B_{\rm pf})$. Substituting $O_{B_{\rm pf}}$ for every
oracle gate and negating the output gives the outer verifier
\[
 D_B(x)=\neg\,C^{O_{B_{\rm pf}}}_{\langle B_{\rm src}\rangle_\nu}
 (x_1,\ldots,x_{\ell_{\rm out}}),\qquad x\in\{0,1\}^q,
\]
which ignores its last $q-\ell_{\rm out}$ inputs. The negation turns
the convention of \cref{imp:transfer-short-pcp} into the usual one: $D_B(x)=1$
means that the outer verifier accepts the random string $x$. Since
$L_\nu\le q^{G'}$, $D_B$ is a circuit of size $q^{O_{{G^\star},C_{\rm clk}}(1)}$ on $q$
inputs. The exponent depends on $C_{\rm clk}$ as well, through
$\mathcal V_{{G^\star}}$; the constants that depend on
it, beginning with the exponent $d_{{G^\star}}$ below, are chosen after
$C_{\rm clk}$ (\cref{earlier:g:app:multitape}). Apply \cref{earlier:g:imp:three-pcpp}. Choose a fixed $d_{{G^\star}}$ bounding
all resulting clause counts and repeat the whole clause list to
exactly
\[
 M=2^{\lceil d_{{G^\star}}\log_2(q+2)\rceil}
\]
clauses. If the list is longer than $M$, which happens at only
finitely many lengths, the machine rejects. The original length divides
this number. This fixes the eventual table arity independently of the
guessed pair.


\subsection{Multitape readings of the refuter and the succinct verifier}\label{earlier:r:app:multitape}\label{earlier:g:app:multitape}\label{earlier:gate-models}

This is \cite[Appendix~C]{NagCG}; its proof is reproduced here, with notation adapted. The Lean companion derives the refuter separately, through its source adapters; the written argument is not a transcription of that formal proof. This appendix supplies the machine-model facts used by the lower-bound
transfer (\cref{app:transfer}). The transfer applies \cref{imp:transfer-refuter} and \cref{imp:transfer-short-pcp} to
multitape machines, the model of \cref{sec:prelim}. The refuter's
proof in \cite[Section~4]{CLW20}, however, assumes random-access
machines for simplicity. The appendix has three parts.
\begin{itemize}
\item \Cref{earlier:r:lem:multitape-refuter} derives the multitape reading of the
refuter from the proof in \cite[Section~4]{CLW20}. It imports the
accepting-time universal simulation of~\cite{SFM78,BGW70}, used with
genuine input and witness prefixes, and it gives
the hierarchy language \(L_{\Tclk}\) a multitape verifier
\(\mathcal D\). \Cref{earlier:r:cor:framed-refuter} then proves the exact framed
convention of the Lean interface, including at linear clocks.
\item \Cref{earlier:r:lem:succinct-verifier} constructs the one fixed machine to
which the short PCP is applied, the succinct hierarchy verifier
\(\mathcal V_G\). It reads a source input through the short description
of a circuit and runs \(\mathcal D\) on the described word.
\Cref{earlier:r:app:short-pcp-proof} proves the short PCP for it.
\item The paragraph after \cref{earlier:r:lem:succinct-verifier} fixes the order
in which the constants of the transfer are chosen.
\end{itemize}

\begin{lemma}[Multitape refuter]\label{earlier:r:lem:multitape-refuter}\label{earlier:g:lem:multitape-refuter}
\Cref{imp:transfer-refuter} holds with every machine in it, the class
\(\NTIME\) and the procedure \(R_{\Tclk}\) read for multitape Turing
machines, with the same quantifier order. The language \(L_{\Tclk}\) has a
nondeterministic multitape verifier that halts within \(O(\Tclk(N))\)
steps on every input and every guess.
\end{lemma}

\emph{Idea of proof.} The diagonal language of Chen, Lyu, and Williams
runs each enumerated machine through a universal simulation. We use the
accepting-time simulation of Seiferas, Fischer, and Meyer for multitape
machines, under an explicit clock of \(O(\Tclk(N))\) steps (Step~1).
For each admitted machine, every verdict required by the diagonal
recurrence has a faithful accepting simulation that survives this
clock at large lengths; other branches may time out. Thus
\cite[Lemma~4.3]{CLW20} holds for it (Step~2). The rest
of their proof uses the machine model only through that lemma and
through SAT-oracle queries, which the Cook--Levin tableau handles for
multitape machines (Step~3). So their refuter applies verbatim
(Step~4), and the guess budgets are unchanged (Step~5).

\begin{proof}
\emph{Conventions.} A nondeterministic multitape machine reads its
guessed bits from a witness tape and is deterministic once the witness
is fixed. It guesses at most \(g(N)\) bits when every run on an input of
length \(N\) reads at most \(g(N)\) witness cells. Fix \(\Tclk\), which is
computable from \(1^N\) in \(O(\Tclk(N))\) steps, and an enumeration of
these machines. As in \cite[Section~4.1]{CLW20}, \(V_M(x,T,z)\) is the
outcome of running \(M\) on \(x\) with witness \(z\) for at most \(T\)
steps. It is \(\bot\) when the run does not halt within \(T\) steps or
reads more than \(\abs z\) witness bits. A machine that guesses at most
\(g(N)\le\abs z\) bits behaves the same with the witness \(z\) and with
its prefix of length \(g(N)\). So it may equally be given exactly
\(g(N)\) witness bits. This proof uses the ordinary witness-tape
convention. \Cref{earlier:r:cor:framed-refuter} treats the framed convention
directly, without adding preprocessing to the weak machine.

\emph{The universal simulation.} We use the accepted-input time
definition of~\cite[p.~147]{SFM78}, its tape-reduction Lemma~3
(pp.~149--150), and its universal-simulation Condition~2 (p.~152).
Composing the tape reduction with that simulator gives a fixed
universal nondeterministic acceptor \(U\) that preserves the language
of each simulated machine \(P\): if \(P\)
has an accepting computation of length \(t\) on \(u\), then \(U\)
on its code and \(u\) has one of length at most \(c_P(t+1)\).
Here \(c_P\) depends only on \(P\); the additive one covers
initialization and zero-step computations. No halting or time-bound
hypothesis is imposed on other inputs or branches. The tape reduction
is due to Book, Greibach, and Wegbreit~\cite{BGW70};
\cite{SFM78} sketches its construction and states Condition~2 as a
property of the program coding. Fortnow and Santhanam use the
linear-time simulation in their hierarchy argument
\cite[proof of Theorem~7.1]{FS16}.

\emph{Shared prefix simulation.} First prepare the actual input and
supplied witness tapes, with their prescribed boundaries and heads at
zero, in a uniform \(c_0N\) setup independent of \(M\). For each
simulation the hierarchy verifier guesses a positive length \(r\) in
unary and copies cells \(0,\ldots,r-1\) of both prepared tapes,
including genuine blank cells when necessary. It packs these two
prefixes, their boundary flags, and cutoff markers at cell \(r\)
into one binary word of length \(O(r+1)\), using a fixed encoding
independent of \(M\). The contents and flags are copied, not guessed.
Two fixed acceptors \(P_M^+\) and \(P_M^-\) unpack this word into
separate tapes, copying the second prefix to a work tape and rewinding
in \(O_M(r+1)\) time. They then execute the deterministic computation
of \(M\) with this witness fixed, accepting only upon, respectively,
an accepting or a rejecting halt of \(M\). Reaching a cutoff aborts
without either verdict, before executing an original transition there.
The cutoff and boundary flags are kept on protected tracks, so writes
by \(M\) cannot erase them. In the present ordinary-witness convention,
an attempt to read beyond the supplied witness has exactly the
\(V_M\) failure behavior \(\bot\), and neither acceptor accepts it.
In the framed convention of \cref{earlier:r:cor:framed-refuter}, the copied
tapes instead include the prescribed frame terminators and blank
zeros beyond them, with no artificial failure at a frame terminator.
These are separate acceptors for the two verdicts; we never complement
the universal acceptor.

For every guessed \(r\), the initial simulated cells agree with the
original tapes. Induction on the simulated steps, including their
writes, preserves the state, heads, and all visited cells until a
cutoff or a genuine halt. Hence every verified verdict is genuine.
Conversely, if the original computation halts in \(t\) steps, the
choice \(r=t+2\) keeps both heads strictly before the cutoffs: a head
starts at zero and moves by at most one cell per step. Prefix copying,
unpacking, and this computation cost \(O_M(t+1)\), and the universal
simulation therefore has a corresponding accepting branch within
\(a_M+b_M(t+1)\) steps, after increasing fixed constants. The choice
of \(r\) is existential in this argument; the algorithm does not
compute \(t\) or a runtime bound for \(M\). The original work tapes
remain unbounded. A computation on some other prefix can run
arbitrarily long without reaching a cutoff, which is why we use an
accepting-time simulation rather than assert a time bound for every
input of \(P_M^+\) or \(P_M^-\).

\emph{Step 1: the hierarchy verifier.} Let \(\mathcal D\) be the following
clocked variant of the algorithm \(A^{\Tclk}_{\rm FS}\) of
\cite[Definition~4.2]{CLW20}. First \(\mathcal D\) computes \(\Tclk(N)\)
and prepares its clock in \(O(\Tclk(N)+1)\) steps.
It then parses its input as \(\langle M,N,z\rangle\) with
\(\abs z=\floor{N/10}\), rejecting otherwise. This takes at most \(c_0N\)
steps for a constant \(c_0\) that does not depend on \(M\), enlarged
to include preparation of the actual tapes described below. It then
uses the shared prefix construction for two universal simulations:
\begin{enumerate}[label=(\alph*)]
\item the rejecting-verdict acceptor \(P_M^-\), for \(M\) on input
\(\langle M,N,0^{\floor{N/10}}\rangle\) with witness \(z\);
\item when \(z\ne1^{\floor{N/10}}\), the accepting-verdict acceptor
\(P_M^+\), for \(M\) on input
\(\langle M,N,z+1\rangle\) with a guessed witness of \(\floor{N/10}\)
bits.
\end{enumerate}
It accepts exactly when the first simulation and, when required, the
second simulation accept. The original run in (a) has a fixed witness
and is deterministic; an accepted first simulation certifies its
rejecting halt, never a cutoff or a failed simulation. The clock
rejects every branch whose steps after clock preparation exceed
\((c_0+2)(\Tclk(N)+1)\), including overlong prefix guesses,
machine-description preparation, and universal-simulation branches.
Hence every branch halts within \(O(\Tclk(N)+1)\) steps on every
input, and \(L_{\Tclk}=L(\mathcal D)\) has the stated verifier.

The language \(L(\mathcal D)\) may differ from \(L(A^{\Tclk}_{\rm FS})\)
on inputs where a simulation exceeds the clock or the encoded machine
breaks its promise. Step~2 shows that it obeys the formula of
\cite[Lemma~4.3]{CLW20} on every input that the proofs use.

\emph{Step 2: Lemma~4.3 of \cite{CLW20} holds for every admitted machine.}
Fix \(M\) that runs in time \(t_M(N)=o(\Tclk(N))\) on every input and
witness and guesses at most \(\floor{N/10}\) bits. Consider an input
\(x^{(N)}_i=\langle M,N,w^{(N)}_i\rangle\) of \cite[Section~4.2]{CLW20}.
Whenever the original computations give the requested verdicts, the
shared prefix argument supplies a branch verifying them with
\(r=t_M(N)+2\). After computing \(\Tclk(N)\) and preparing the
clock, its total cost is at most
\[
 c_0N+a_M+2b_M(t_M(N)+1)
 <(c_0+2)(\Tclk(N)+1)
\]
for all sufficiently large \(N\), since \(N\le\Tclk(N)\) and
\(t_M(N)=o(\Tclk(N))\). The constants \(a_M,b_M\) cover the fixed
machine descriptions and prefix simulations; only \(c_0\) is uniform
in \(M\). Thus every required verdict has an accepting simulation
that survives the cap, whereas faithfulness rules out any branch
certifying a false verdict. Other branches can time out. Every
run of \(M\) halts within \(t_M(N)\le\Tclk(N)\) steps and reads at most
\(\floor{N/10}\) witness bits. Hence \(\mathcal D\) accepts
\(x^{(N)}_i\) exactly when \(V_M(x^{(N)}_1,\Tclk(N),w^{(N)}_i)=0\) and,
for \(i<2^{\floor{N/10}}\), \(M(x^{(N)}_{i+1})=1\). This is the formula
of \cite[Lemma~4.3]{CLW20}, which therefore holds for \(\mathcal D\) and
\(M\) at every sufficiently large \(N\). The overhead is the constant
factor \(b_M\), so the little-\(o\) promise needs no logarithmic slack.

\emph{Step 3: the oracle queries.} The proofs of \cite[Theorems~1.12
and~4.6]{CLW20} use the machine model only through Lemma~4.3 and through
SAT-oracle queries of three kinds:
\begin{itemize}
\item whether \(\mathcal D\) accepts a given input;
\item whether \(M\) accepts a given input within \(\Tclk(N)\) steps with
some witness in a given lexicographic interval;
\item whether \(V_M(x,\Tclk(N),w)=1\) for given \(x\) and \(w\).
\end{itemize}
Each query concerns runs of \(O(\Tclk(N))\) steps of one nondeterministic
multitape machine. The Cook--Levin tableau turns it into a CNF formula
of size \(O(\Tclk(N)^2)\). The formula is constructed in
\(\poly(\Tclk(N))\) time by a polynomial whose degree does not depend on
\(M\); for queries about \(M\), only the constant factor depends on
\(M\).
For the companion's exact SAT oracle, apply the standard equisatisfiable
Tseitin conversion and repeat literals in unit and binary clauses so
that every clause has exactly three literals; an empty clause is
represented by a fixed contradictory pair of padded unit clauses.
Encode each such clause using the companion's native literal-list
coding, and encode the formula using its balanced clause-list coding.
Because the clause length is fixed and the outer list is balanced,
these encodings and the query framing have polynomial length and
construction time, with a degree independent of \(M\). Thus the oracle
answers and the stated polynomial bounds are preserved.

\emph{Step 4: the refuter.} With Steps~2 and~3, the proof of
\cite[Theorem~4.6]{CLW20} applies verbatim. \(R_{\Tclk}\) first asks
whether \(M\) accepts \(x^{(N)}_1\). If not, it runs the binary search of
\cite[Corollary~4.5]{CLW20} on the lists \(\mathcal D(x^{(N)}_i)\) and
\(M(x^{(N)}_i)\), \(1\le i\le2^{\floor{N/10}}\). If so, it finds by
binary search the least \(j\) with \(V_M(x^{(N)}_1,\Tclk(N),w^{(N)}_j)=1\)
and runs the same search on the first \(j\) entries. Either way it makes
\(O(N)\) adaptive queries and runs in \(\poly(\Tclk(N))\) time, with a
degree that does not depend on \(M\). At every sufficiently large \(N\) it
returns a word \(x^{(N)}_p\) with \(\mathcal D(x^{(N)}_p)\ne M(x^{(N)}_p)\).

\emph{Step 5: the guess budgets.} The guesses made by the simulations in
Step~1 are guesses of \(\mathcal D\), not of \(M\), so the promise on \(M\)
is unchanged. The weak machine of
\cref{app:transfer} guesses at most
\(\floor{N/16}\le\floor{N/10}\) bits and runs in time \(o(\Tclk(N))\), so
the lemma applies to it.
\end{proof}

\begin{corollary}[Refuter for framed machines]\label{earlier:r:cor:framed-refuter}\label{earlier:g:cor:framed-refuter}
The conclusion of \cref{imp:transfer-refuter} also holds with the following
convention, for every time-constructible
\(N\le\Tclk(N)\le2^{\poly(N)}\). Set \(g(N)=\floor{N/16}\) and
\[
 \operatorname{fr}(b_1\cdots b_m)=1b_1\,1b_2\cdots1b_m\,0,
 \qquad \operatorname{fr}(\epsilon)=0.
\]
A weak machine is a deterministic multitape verifier whose initial
tapes contain \(\operatorname{fr}(x)\) and
\(\operatorname{fr}(w)\), with \(\abs x=N\) and \(\abs w=g(N)\);
all other tapes are blank, all heads start at the leftmost cell, and
blank cells contain zero. It halts within \(t_M(N)=o(\Tclk(N))\)
steps for every such \(x,w\), and \(M(x)=1\) means that some such
\(w\) is accepted. The hierarchy language has a verifier with the
same framing convention, time bound \(c(\Tclk(N)+1)\), and exactly
\(c(\Tclk(N)+1)\) witness bits, for a fixed positive integer \(c\).
The hierarchy language and the refuter are fixed before \(M\); only
the eventual onset and the coefficient in the refuter's polynomial
time bound may depend on \(M\).
\end{corollary}

\begin{proof}
We apply the diagonal construction in the proof of
\cref{earlier:r:lem:multitape-refuter} directly to an enumeration of these
framed verifiers. Thus \(M\) retains its original runtime and its
original initial tapes. The two conventions may give different
hierarchy languages, as \cref{imp:transfer-refuter} permits.

For fixed \(M,N\), enumerate all \(g(N)\)-bit strings as
\(w_1,\ldots,w_s\), where \(s=2^{g(N)}\), and use the length-\(N\)
records \(x_i=\langle M,N,w_i\rangle\) of~\cite[Section~4.1]{CLW20}.
For every fixed \(M\) these records exist at all sufficiently large
\(N\). Write \(V_M^{\rm fr}(x,t,w)\) for the outcome of the original
verifier within \(t\) steps on exactly the two framed tapes above,
with value \(\bot\) if it has not halted. In particular the terminator
is placed immediately after the \(g(N)\)th witness bit; no longer
witness is substituted for \(w\).

\emph{Uniform setup cost.} On the record \(x_i\), the hierarchy
verifier prepares \(\operatorname{fr}(x_1)\) and
\(\operatorname{fr}(w_i)\) for the first simulation. For the second
simulation, when \(i<s\), it prepares \(\operatorname{fr}(x_{i+1})\)
and \(\operatorname{fr}(v)\), where it guesses exactly \(g(N)\) bits
\(v\). Parsing the outer record, incrementing its witness field,
writing these frames, and returning their heads to zero cost at most
\(c_0N\) steps, for an absolute constant \(c_0\). The frames have
lengths \(2N+1\) and \(2g(N)+1\). Their bit representation is
independent of \(M\); it does not encode every data bit by a
machine-dependent block.

Apply the shared prefix simulation from the proof of
\cref{earlier:r:lem:multitape-refuter} to these prepared framed tapes. The
copied data include the actual terminator immediately after the
prescribed witness, if it lies within the prefix, and actual blank
zeros beyond it. Thus the deterministic computations being verified
have exactly the initial tapes in the statement, even if \(M\)
reads or writes past a frame terminator. Only the separate protected
prefix cutoff aborts without a verdict. In particular an aborted
prefix does not supply the rejecting verdict required in the first
simulation.

For each promised computation choose \(r=t_M(N)+2\) in that argument.
The fixed-machine decoding cost is \(a_M\), and prefix processing
and a useful accepting simulation cost at most
\(b_M(t_M(N)+1)\) per run. Whenever the original computations give
the requested verdicts, there is therefore a branch verifying both
with total cost at most
\[
 c_0N+a_M+2b_M(t_M(N)+1).
\]
Here \(c_0\) is independent of the enumerated machine, while
\(a_M,b_M\) may depend on it. In particular there is no
machine-dependent coefficient multiplying the initial data length.

After computing \(\Tclk(N)\) and preparing a clock, impose the
uniform step cap \((c_0+2)(\Tclk(N)+1)\), as in Step~1 of
\cref{earlier:r:lem:multitape-refuter}. For each fixed promised \(M\),
\[
 c_0N+a_M+2b_M(t_M(N)+1)
 <(c_0+2)(\Tclk(N)+1)
\]
eventually, since \(N\le\Tclk(N)\) and
\(t_M(N)=o(\Tclk(N))\). The cap rejects overlong branches; an
invalid simulation cannot certify a verdict. Every required verdict
has an accepting simulation within the cap at these lengths, and
faithfulness prevents spurious accepting branches.
The resulting diagonal language, say \(L^{\rm fr}_{\Tclk}\), thus
satisfies the exact recurrence
\[
 L^{\rm fr}_{\Tclk}(x_i)=
 \ind[V_M^{\rm fr}(x_1,\Tclk(N),w_i)=0]\cdot
 \begin{cases}
 M(x_{i+1}),&i<s,\\
 1,&i=s.
 \end{cases}
\]
This is the identity used by~\cite[Lemma~4.3 and Theorem~4.6]{CLW20},
with \(g(N)\) in place of \(\floor{N/10}\). Its proof and the
binary-search refuter use only the complete enumeration of witnesses,
so Steps~3--4 of \cref{earlier:r:lem:multitape-refuter} apply. The SAT queries
now encode the specified framed tapes and their fixed terminators.
Their initial descriptions have length \(O(N)\), and their time
bounds are \(O(\Tclk(N))\); the same polynomial tableau bounds and
\(O(N)\) query count follow. Encoding the refuter's request and output
with the companion's nested frame encoding changes these polynomial bounds
only by constant factors. The polynomial degree is fixed with the
refuter, before \(M\).

Finally the clocked diagonal verifier uses \(O(\Tclk(N)+1)\) steps
and nondeterministic bits on every branch. A deterministic verifier
reads those choices from a framed witness, decoding each requested bit
in constant time. Taking a sufficiently large fixed integer \(c\)
gives both the step bound and the witness length
\(c(\Tclk(N)+1)\); unused witness bits are ignored. This is the
hierarchy convention in the statement. All setup and simulation
guesses belong to this hierarchy verifier. The weak machine still has
exactly \(g(N)\) witness bits and its original runtime. The argument
uses only \(N\le\Tclk(N)\), so it includes linear clocks.
\end{proof}

\begin{lemma}[Succinct hierarchy verifier]\label{earlier:r:lem:succinct-verifier}\label{earlier:g:lem:succinct-verifier}
Fix integers \(G\ge1\), \(G'\ge2\) and \(C_{\rm clk}\ge2\), let
\(\Tclk(N)=N^{C_{\rm clk}}\), let \(\mathcal D\) be the verifier of
\(L_{\Tclk}\) from \cref{earlier:r:lem:multitape-refuter}, and let
\(L_\nu=\nu^{G'}\). Use the self-delimiting binary description of
general circuits fixed in \cref{app:transfer}, and let
\(\mathrm{tt}_N\) and \(\langle B\rangle_\nu\) be as there. There are a deterministic multitape machine
\(\mathcal V_G(u,y)\), with input \(u\) and witness \(y\), and a constant
\(c_G\ge1\) such that, with
\[
 T_G(n)=c_G\,2^{C_{\rm clk}\floor{n^{1/G'}}},
\]
the following hold.
\begin{enumerate}[label=(\roman*)]
\item \(T_G\) is nondecreasing, \(T_G(n)\ge n\) for every \(n\ge1\),
\(T_G(n)\) is computable from \(n\) in \(\poly(n)\) steps, and
\(T_G(L_\nu)=c_G\Tclk(2^\nu)\) for every \(\nu\ge1\).
\item For every \(n\ge1\), every \(u\in\{0,1\}^n\) and every \(y\),
\(\mathcal V_G\) halts within \(T_G(n)\) steps.
\item Let \(\nu\ge1\) and \(N=2^\nu\), and let \(B\) be a general circuit
with at least \(\nu\) inputs whose description has at most \(L_\nu\)
bits. Some \(y\) of length \(T_G(L_\nu)\) makes
\(\mathcal V_G(\langle B\rangle_\nu,y)\) accept if and only if
\(\mathrm{tt}_N(B)\in L_{\Tclk}\).
\end{enumerate}
\end{lemma}

\emph{Idea of proof.} The verifier \(\mathcal V_G\) checks that its
input has the length \(L_\nu\) of a padded description. It then expands
the described circuit \(B\) into its truth table \(\mathrm{tt}_N(B)\)
and runs \(\mathcal D\) on that word with its own witness. Expanding
takes \(O(NL_\nu^3)\) steps and running \(\mathcal D\) takes
\(O(\Tclk(N))\) steps, so at the lengths \(L_\nu\) the time bound
\(T_G\) is a constant multiple of \(\Tclk(N)\). At every other length
\(\mathcal V_G\) rejects within \(O(n)\) steps.

\begin{proof}
On \((u,y)\) with \(\abs u=n\), the machine \(\mathcal V_G\) works in four
stages.
\begin{enumerate}[label=(\arabic*)]
\item It writes \(n\) in binary, computes \(\nu=\floor{n^{1/G'}}\), and
rejects unless \(\nu\ge1\) and \(n=\nu^{G'}\). This takes \(O(n)\)
steps.
\item It parses \(u\) as the description of a general circuit \(B\)
followed by zeros. It rejects if the description is malformed or \(B\)
has fewer than \(\nu\) inputs. This takes \(O(n)\) steps.
\item For \(i=0,1,\ldots,N-1\), where \(N=2^\nu\), it evaluates \(B\) at
\(\mathrm{bin}_\nu(i)0^{a-\nu}\), where \(a\) is the number of inputs of
\(B\). It writes the value into cell \(i\) of a work tape, which then
holds \(\mathrm{tt}_N(B)\). The circuit has at most \(n\) gates. Keeping
the gate values on a tape and fetching the two arguments of each gate by
a scan evaluates it in \(O(n^3)\) steps. So this stage takes
\(O(Nn^3)\) steps.
\item It runs \(\mathcal D\) on \(\mathrm{tt}_N(B)\) and accepts exactly
when \(\mathcal D\) accepts. The work tape of stage~(3) serves as the
input tape of \(\mathcal D\), its own witness tape as the witness tape
of \(\mathcal D\), and further work tapes as those of \(\mathcal D\). By
\cref{earlier:r:lem:multitape-refuter} this takes at most
\(c_{\mathcal D}\Tclk(N)\) steps for a constant \(c_{\mathcal D}\).
\end{enumerate}
If \(n=L_\nu\), the total is
\[
 O(n)+O(NL_\nu^3)+c_{\mathcal D}\Tclk(N)\le c'2^{C_{\rm clk}\nu}
\]
for a constant \(c'\), because \(L_\nu^3=\nu^{3G'}\) and
\(C_{\rm clk}\ge2\). Otherwise the run stops in stage~(1) within
\(c''n\) steps, and \(n<(\nu+1)^{G'}\) with \(\nu=\floor{n^{1/G'}}\).
Choose \(c_G\ge c'\) so large that
\(c_G2^{C_{\rm clk}\nu}\ge(c''+1)(\nu+1)^{G'}\) for every \(\nu\ge1\).
This is possible because \((\nu+1)^{G'}2^{-C_{\rm clk}\nu}\) is bounded.
Then (ii) holds, and so does \(T_G(n)\ge n\). The other assertions of (i)
are immediate from the definition.

For (iii), if \(\mathcal V_G(\langle B\rangle_\nu,y)\) accepts, then
\(\mathcal D\) accepts \(\mathrm{tt}_N(B)\) with a witness, so
\(\mathrm{tt}_N(B)\in L(\mathcal D)=L_{\Tclk}\). Conversely, suppose that
\(\mathrm{tt}_N(B)\in L_{\Tclk}\), and take an accepting run of
\(\mathcal D\) on it. The run has at most
\(c_{\mathcal D}\Tclk(N)\le T_G(L_\nu)\) steps, so it reads only witness
cells below \(T_G(L_\nu)\). Writing the witness bits it reads into a word
\(y\) of length \(T_G(L_\nu)\), with zeros elsewhere, gives the same run,
and \(\mathcal V_G(\langle B\rangle_\nu,y)\) accepts.
\end{proof}

\paragraph{Use in the transfer and the order of constants.} In
\cref{app:transfer}, \(\mathcal D\) is the
multitape verifier of \(L_{\Tclk}\) from \cref{earlier:r:lem:multitape-refuter},
with \(\Tclk(N)=N^{C_{\rm clk}}\), and \cref{imp:transfer-short-pcp} is applied to
\(\mathcal V_G\) with the time bound \(T_G\). The constants are chosen
in this order.
\begin{enumerate}
\item \(d_0=2\). It
bounds the exponent of \(N\) in the check \eqref{earlier:g:eq:source-check} and
depends on no PCP construction.
\item The general-circuit exponent \(G\), from \(h_{\rm sim}\) and
\(b_{\rm amp}\); it does not depend on \(C_{\rm clk}\).
\item The exponent \(G'\ge2\), from \(G\) and the description format of
general circuits (\cref{app:transfer}). It fixes
\(L_\nu=\nu^{G'}\), and it does not depend on \(C_{\rm clk}\) or
\(A_{\rm pad}\).
\item An integer \(C_{\rm clk}>d_0+2\). This fixes \(\mathcal D\), and
\cref{earlier:r:lem:succinct-verifier} then fixes \(\mathcal V_G\), \(c_G\) and
\(T_G\).
\item \Cref{imp:transfer-short-pcp}, applied to \(\mathcal V_G\) and \(T_G\), fixes
a constant \(c_\ell\ge1\) with
\[
 \ell\le\lgx T_G(n)+c_\ell\lgx(\lgx T_G(n)+2)
\]
for every \(n\), and the polynomials bounding the size of \(C^O_u\) and
the running time of its construction. (\Cref{earlier:r:thm:short-pcp-ours} derives
them from the constants of~\cite{BCGT13} and of
\cref{earlier:r:lem:acsp-degree,earlier:r:lem:generator}.)
\item An integer \(A_{\rm pad}\ge c_\ell\) with
\(A_{\rm pad}\ge b+c_\ell\lgx(b+1)\), where \(b=\lgx c_G\). For
\(N=2^\nu\) write \(\tau=\lgx\Tclk(N)=C_{\rm clk}\nu\). Since
\(T_G(L_\nu)=c_G2^\tau\) and \(\tau+b+2\le(\tau+2)(b+1)\),
\[
 \ell_{\rm out}\le\tau+b+c_\ell\lgx(\tau+b+2)
 \le\tau+c_\ell\lgx(\tau+2)+b+c_\ell\lgx(b+1)\le q(N).
\]
So the outer verifier fits the native width. Its oracle queries
\(a\in\{0,1\}^{\ell_{\rm out}}\) can be zero-padded to the \(q\)
inputs of \(B_{\rm pf}\), and the tables \(X_\zeta\) and
\(\widehat O^*\) of \cref{earlier:g:sec:gadget-extract} have
\(q\) inputs. Moreover
\[
 2^{q(N)}\le2\cdot4^{A_{\rm pad}}\Tclk(N)(\lgx\Tclk(N)+2)^{A_{\rm pad}},
\]
so \(2^q/\Tclk(N)=q^{O(1)}\).
\item Then the remaining constants, in this order:
\begin{itemize}
\item the clause-count exponent \(d_G\);
\item finally the padding and cap constants: \(K_{\rm amp}\),
\(B_{\rm len}\) and \(K_d\) in
\cref{app:transfer}.
\end{itemize}
The constant \(\eta_0\) of \cref{app:transfer} depends only on
\cref{thm:counter}.
\end{enumerate}
No constant depends on one chosen after it. In particular,
\(\mathcal V_G\) bounds the size of a source circuit only through the
length \(L_\nu\). That length depends on \(N\) and \(G'\) alone, not on
\(q\) or \(A_{\rm pad}\). The weak machine enforces the bound \(q^G\)
itself.

\subsection{Extracting the weighted hard table}\label{earlier:gate-canonical}

This reproduces \cite[Section~11.5]{NagCG}, with notation adapted. \label{earlier:g:sec:gadget-extract}

Fix such a rejected yes-instance $\zeta\in\{0,1\}^N$, with $N=2^\nu$. The
two cases below, the
SAT-checked canonical choices, and the clause-indexed proof table
follow Chen--Lyu--Williams~\cite[Sections~6.2.4--6.2.5]{CLW20}.
Their proof of~\cite[Lemma~6.5]{CLW20} averages cofactors into
real-valued proofs; we instead select one
Boolean cofactor per name and weight the list positions. Call a general
circuit with $q$ inputs and at most $q^{G^\star}$ gates a \emph{source circuit
for $\zeta$} if it satisfies \eqref{earlier:g:eq:source-check}. Descriptions are
ordered lexicographically, and ``first'' refers to this order. The
canonical choices are made in three steps, which yield the three cases
named in \cref{sec:ae-extract}.
\begin{enumerate}[label=(\roman*)]
\item (\textsc{NoSource}) If $\zeta$ has no source circuit, let $X_\zeta:\{0,1\}^q\to\{0,1\}$ map
$u$ to $\zeta_i$, where $\mathrm{bin}_\nu(i)$ is the first $\nu$ coordinates
of $u$. A general circuit of at most $q^{G^\star}$ gates computing $X_\zeta$ would
be a source circuit for $\zeta$. Hence $X_\zeta$ has general circuit complexity
greater than $q^{G^\star}$.
\item Otherwise let $B^*_{\rm src}$ be the first source circuit for $\zeta$.
Then $\mathrm{tt}_N(B^*_{\rm src})=\zeta\in L_{\Tclk}$, so some witness of
length $T_{{G^\star}}(L_\nu)$ makes $\mathcal V_{{G^\star}}$ accept
$\langle B^*_{\rm src}\rangle_\nu$, and \cref{imp:transfer-short-pcp} provides an
oracle on $\ell_{\rm out}$ inputs that makes
$C^O_{\langle B^*_{\rm src}\rangle_\nu}$ unsatisfiable, that is,
accepted on every random string. Let $O^*$ be the first such oracle,
viewed as a truth table, and let $\widehat O^*(av)=O^*(a)$ for
$a\in\{0,1\}^{\ell_{\rm out}}$ and $v\in\{0,1\}^{q-\ell_{\rm out}}$.
\item (\textsc{NoProof}) If no general circuit $B_{\rm pf}$ with $q$ inputs and at most
$q^{G^\star}$ gates satisfies $O_{B_{\rm pf}}=O^*$, then $\widehat O^*$ has
general circuit complexity greater than $q^{G^\star}$, since a circuit computing
it would satisfy that equation.
\end{enumerate}
Cases (\textsc{NoSource}) and (\textsc{NoProof}) are treated together. In each of them
we obtain a function on $q$ inputs with general circuit complexity
greater than $q^{G^\star}$. By
\cref{imp:transfer-amplification} and the choice of ${G^\star}$, the
amplified function, on at most $a_{{G^\star}}q$ inputs, has agreement $1/2+o(1)$
against every target circuit considered here. This is eventually below
$1-\gamma/5+\epsilon$.

Otherwise, in case (\textsc{Tautology}), let $B^*_{\rm pf}$ be the first
general circuit with $q$ inputs and at most $q^{G^\star}$ gates such that
$O_{B^*_{\rm pf}}=O^*$, and put
$B=(B^*_{\rm src},B^*_{\rm pf})$. Then $D_B$ is a tautology, because
$C^{O^*}_{\langle B^*_{\rm src}\rangle_\nu}$ is unsatisfiable, and its
source circuit passes the check \eqref{earlier:g:eq:source-check}, so the weak
machine can guess $B$. This canonical choice suffices: the argument
below uses only that $B$ is a pair whose source circuit passes the
check, that $D_B$ is a tautology, and that the weak machine rejected $\zeta$
on every branch, including those that guess $B$. The efficient honest proof from
\cref{earlier:g:imp:three-pcpp} defines Boolean coordinate functions $\Pi_X(x)$
that make every three-literal clause true. For each fresh helper, define its
honest function to be the AND of that clause's three honest literal
values. For gadget $a$, form the five-entry list of underlying
coordinate names
\[
 (X_{a,1},X_{a,2},X_{a,3},Y_a,Y_a).
\]
The first three entries are unsigned names: the clause's literal
signs remain in the verifier, not in the table definition. Fill
$P_q$ address positions by cycling through this list, and define
the target table to return the corresponding honest coordinate.
Its exact arity is
\begin{equation}\label{earlier:g:eq:gadget-table-arity}
 q_0(q)=q+\log_2M+\log_2P_q=q+O(\log q).
\end{equation}



\paragraph{Weighted-cofactor incidence.}\label{earlier:weighted-incidence}
Let \(M_{\rm cl}=M\) be the number of gadgets, let \(P_q\ge5\) be the positions per gadget, and put \(t=\lfloor P_q/5\rfloor\).
Suppose a Boolean candidate table \(C(x,\alpha)\) has uniform error \(\delta\) against the honest table above. For each auxiliary unsigned name \(X\), choose an address cofactor minimizing its error \(e_X\) against \(\Pi_X\); systematic names use their prescribed parities and have \(e_X=0\).
Let \(I_X\) be its total stability weight across the gadgets: one for each literal position and two for each helper. The five-entry lists repeat every name in at least \(tI_X\) table positions, so
\[
 t\sum_X I_Xe_X\le M_{\rm cl}P_q\delta.
\]
Weighted gadget stability (\cref{earlier:g:lem:gadget-stability}), averaged over gadgets and inputs, therefore gives
\[
 \E Q\le\frac1{M_{\rm cl}}\sum_X I_Xe_X
 \le\frac{P_q}{t}\delta\le(5+\epsilon)\delta.
\]
For the paper's choice \(P_q=128/\epsilon\), with \(0<\epsilon<1\), we have \(t\ge4/\epsilon\), so \(P_q/t\le5+4/t\le5+\epsilon\).
Repeated names are counted with their multiplicities, complements preserve error, and systematic positions only strengthen the table-error bound. This is the incidence argument of~\cite[Equation~(18)]{NagCG}.

\subsection{Uniform recovery of the canonical choices}\label{earlier:uniform-recovery}
This is \cite[Section~11.6]{NagCG}; the proof is reproduced with
notation adapted. 

All canonical choices are recoverable in $\ENP$.
The fixed-machine refuter runs in $\operatorname{poly}(\Tclk(N_\nu))
=2^{O(q)}$ time with SAT access. The choices of
\cref{app:transfer} are recovered in three stages, by prefix
self-reduction on SAT formulas of size $2^{O(q)}$. First, a formula over
the description of $B_{\rm src}$ with the $N$ constraints
\eqref{earlier:g:eq:source-check}, of size $Nq^{O({G^\star})}$, finds the first source
circuit $B^*_{\rm src}$, or shows that there is none, in which case the
table is $X_\zeta$. Second, the algorithm of \cref{imp:transfer-short-pcp} is run
deterministically on $\langle B^*_{\rm src}\rangle_\nu$, in
$q^{O_{{G^\star},C_{\rm clk}}(1)}$ time; a formula with one variable per oracle address and the
constraints $C^O_{\langle B^*_{\rm src}\rangle_\nu}(R)=0$ for all
$2^{\ell_{\rm out}}$ inputs $R$, of size
$2^{2\ell_{\rm out}}q^{O_{{G^\star},C_{\rm clk}}(1)}=2^{O(q)}$ (each oracle gate
becomes a selector over the $2^{\ell_{\rm out}}$ oracle variables),
finds the first oracle $O^*$. Third, a formula over the description of
$B_{\rm pf}$ with the $2^{\ell_{\rm out}}$ constraints
$B_{\rm pf}(a0^{q-\ell_{\rm out}})=O^*(a)$, of size
$2^{\ell_{\rm out}}q^{O({G^\star})}$, finds the first $B^*_{\rm pf}$, or shows
that there is none, in which case the table is $\widehat O^*$. No formula
contains the algorithm of \cref{imp:transfer-short-pcp}. Run the general amplifier
on $X_\zeta$ or $\widehat O^*$. If no oracle $O^*$ exists at an exceptional
initial source length, the second stage detects that case and the slice
is set to zero. In the Tautology case, the honest PCPP proof and each AND helper are
computable in polynomial time once the pair is known. The fixed-error schedule of \cref{sec:ae-params} is computable in polynomial time in $q$. Since $q=O(n)$,
the final functions belong to $\ENP$.


\subsection{Exact conditional cycle potentials}\label{earlier:conditional-potentials}
This section gives the conditional-expectation calculation used by the counter.

\begin{corollary}[Conditional potentials]\label{cor:conditional}
Use the generator of \cref{imp:seed} with \(h=m\) and seed length
\(S\). For a seed prefix \(\pi\) put \(c_\pi(w)=\E[\Gamma_w\mid\pi]\).
Then \(2^Sc_\pi(w)\in\Z\), \(|c_\pi(w)|\le(N^2Q)^{2m}\), \(c_\pi(w)\)
does not depend on \(z\), it is computable exactly in
\((Q+1)^{2m}\poly(n,S)\) time, and
\[
 \E[\Phi_z\mid\pi]=\frac{2|\Theta|}{4^m}\sum_wc_\pi(w)\,\mask_w(z)\sum_\sigma\kappa_\sigma
 \ind[Y_{1,\sigma}(z,w)=0]\,\ind[Y_{2,\sigma}(z,w)=0].
\]
For the empty prefix, \(\E\Phi_z\) equals its value under independent
uniform samples.
\end{corollary}
\begin{proof}
Expand \(\Gamma_w=\prod_{l=1}^{2m}\bigl(Q-N^2\sum_{j=1}^Q\ind[(U_j,V_j)=e_l]\bigr)\)
into \((Q+1)^{2m}\) products, each an integer times a conjunction of at
most \(2m\) events ``sample \(j\) lies at \(e_l\)''. By
\cref{imp:seed} each conjunction has conditional probability \(0\) or
\(2^{-a}\) with \(a\le S\), computed by Gaussian elimination; repeated
sample labels or repeated walk points only repeat or augment the
linear system, and a label required to lie at two different points
gives probability \(0\). Since \(\mask_w(z)\) and the indicators do not
depend on the seed, linearity gives the display. The bound
\(|\Gamma(x)|\le N^2Q\) gives the bound on \(|c_\pi|\). Each
conjunction involves at most \(2m\) distinct labels, and the positions
are \(2m\)-wise independent and uniform, so the unconditional
expectation of every term equals its value under independent samples.
\end{proof}

\subsection{The short PCP for the verifiers used here}\label{earlier:r:app:short-pcp-proof}\label{earlier:release-short-pcp}

This is \cite[Appendix~D]{NagCG}; the complete proof of \cite{NagCG} is reproduced with notation adapted.
Its two published premises are the BCGT verifier and Shoup's
irreducible-polynomial algorithm from \cref{par:five-inputs}. This appendix proves \cref{imp:transfer-short-pcp}, the short PCP in the form of
Murray and Williams, for deterministic multitape verifiers with
computable time bounds (\cref{earlier:r:thm:short-pcp-ours}). This case covers its
only use. The transfer applies \cref{imp:transfer-short-pcp} to one verifier: the
succinct hierarchy verifier \(\mathcal V_G\) of
\cref{earlier:r:lem:succinct-verifier}, in \cref{app:transfer}.
We use the PCP construction of Ben-Sasson, Chiesa, Genkin, and Tromer
for algebraic constraint problems~\cite{BCGT13} and import two
results: their strong Reed--Solomon proximity verifiers, with the
soundness bounds of~\cite[Theorem~11.1]{BCGT13}, and Shoup's deterministic
algorithm for irreducible polynomials~\cite[Corollary~3.2]{Sho88}, as
used in~\cite[Algorithm~17]{BCGT13}. The reduction from the verifier to
an algebraic constraint problem is proved here. References to numbered
statements of~\cite{BCGT13} are to its full version.

\paragraph{Plan of the proof.} The verifier is first simulated by a
random-access machine whose run has an exact length. That run is then
written as an assignment to the points of an affine subspace of a
binary field, so that local polynomial equations hold everywhere on the
subspace exactly when the run accepts. The PCP of~\cite{BCGT13} checks
such an assignment. It amplifies the soundness of its proximity
sub-verifiers by naive sequential repetition, which multiplies the
randomness by a polylogarithmic factor~\cite[Remark~11.2 and
Lemmas~C.6--C.8]{BCGT13}, whereas \cref{imp:transfer-short-pcp} has randomness
\(\lgx T+O(\log\log T)\). We therefore run each sub-verifier once and
amplify the whole verifier with a small-bias set, as
in~\cite[Section~4.4]{BV14}. The appendix has the following parts.
\begin{itemize}
\item \Cref{earlier:r:lem:small-bias,earlier:r:lem:patch-r}: small-bias sets and the
amplification that uses them; \cref{earlier:r:lem:rs-recursion}: the randomness
and proof length of one run of a Reed--Solomon proximity verifier.
\item \Cref{earlier:r:def:ram,earlier:r:lem:tm-ram}: random-access machines, and the
simulation of a multitape verifier by one whose run has an exact
length.
\item \Cref{earlier:r:def:acsp}: affine constraint instances, the problem the PCP
checks.
\item \Cref{earlier:r:thm:trace}: an accepting run as an event trace, checked by
local equations in time order and in address order.
\item \Cref{earlier:r:thm:embedding}: the event trace as an affine constraint
instance, with a routing network joining the two orders;
\cref{earlier:r:lem:acsp-degree,earlier:r:lem:generator}: its degree, its field, and its
uniform construction.
\item \Cref{earlier:r:thm:short-pcp-ours}: assembly in four steps.
\item The routine proof of \cref{earlier:r:lem:tm-ram} comes last.
\end{itemize}

\emph{Notation.} As in \cref{imp:transfer-short-pcp}, \(n\) and \(T(n)\) are the
input length and the time bound of a general verifier. All other
symbols introduced in this appendix are local to it. Logarithms are
base two. Write \(\alpha\cdot\sigma\) for the inner product over
\(\mathbb F_2\).

\subsection{Small-bias sets and one Reed--Solomon test}

\begin{lemma}[Small-bias sets]\label{earlier:r:lem:small-bias}
Let \(r\ge1\) and \(0<\varepsilon<1\). A multiset
\(S\subseteq\{0,1\}^r\) is \emph{\(\varepsilon\)-biased} if
\(\abs{\E_{\sigma\in S}(-1)^{\alpha\cdot\sigma}}\le\varepsilon\) for every
nonzero \(\alpha\in\{0,1\}^r\).
\begin{enumerate}[label=(\roman*)]
\item An \(\varepsilon\)-biased multiset of size less than
\((2r/\varepsilon)^2\) can be computed deterministically in
\(\poly(r/\varepsilon)\) steps.
\item If \(S\) is \(\varepsilon\)-biased and
\(\mathcal B\subseteq\{0,1\}^r\) has density \(\mu>0\), then the set of
\(R\in\{0,1\}^r\) with \(R\oplus\sigma\notin\mathcal B\) for every
\(\sigma\in S\) has density at most \(\varepsilon^2/\mu\).
\end{enumerate}
\end{lemma}
\begin{proof}
(i) This is the powering construction of Alon, Goldreich, H\aa{}stad,
and Peralta~\cite{AGHP92}; see also Naor and Naor~\cite{NN93}. Let \(m\)
be the least integer with \(2^m\ge r/\varepsilon\), so that
\(2^m<2r/\varepsilon\). Find an irreducible polynomial of degree \(m\)
over \(\mathbb F_2\) by exhaustive search, in \(2^m\poly(m)\) steps, and
identify \(\mathbb F_{2^m}\) with \(\mathbb F_2^m\). For
\(a,b\in\mathbb F_{2^m}\) let \(\sigma_{a,b}\in\{0,1\}^r\) have
coordinates \(a^j\cdot b\), \(0\le j<r\), and let \(S\) consist of all
\(2^{2m}<(2r/\varepsilon)^2\) vectors \(\sigma_{a,b}\). For nonzero
\(\alpha\), \(\alpha\cdot\sigma_{a,b}=p_\alpha(a)\cdot b\), where
\(p_\alpha(X)=\sum_j\alpha_jX^j\) is a nonzero polynomial of degree less
than \(r\). If \(p_\alpha(a)\ne0\), the bit \(p_\alpha(a)\cdot b\) is
unbiased over \(b\), and \(p_\alpha(a)=0\) for fewer than \(r\) values
of \(a\). Hence the bias is less than \(r/2^m\le\varepsilon\).

(ii) Let \(g(R)=\E_{\sigma\in S}\ind_{\mathcal B}(R\oplus\sigma)\), and
write \(\widehat h(\alpha)=\E_R h(R)(-1)^{\alpha\cdot R}\). Then
\(\E_Rg(R)=\mu\) and
\(\widehat g(\alpha)=\widehat{\ind_{\mathcal B}}(\alpha)
\E_{\sigma\in S}(-1)^{\alpha\cdot\sigma}\). By Parseval,
\[
 \E_R\bigl(g(R)-\mu\bigr)^2=\sum_{\alpha\ne0}\widehat g(\alpha)^2
 \le\varepsilon^2\sum_{\alpha\ne0}\widehat{\ind_{\mathcal B}}(\alpha)^2
 =\varepsilon^2(\mu-\mu^2).
\]
On the set in question \(g=0\), so \((g-\mu)^2=\mu^2\) there, and its
density is at most \(\varepsilon^2(\mu-\mu^2)/\mu^2\le\varepsilon^2/\mu\).
\end{proof}

Part (ii) is the hitting property~\cite[Lemma~4.7]{BV14}. Equivalently,
it follows from the expander mixing lemma of Alon and
Chung~\cite{AC88} for the Cayley graph of \(\mathbb F_2^r\) generated by
\(S\), whose nontrivial eigenvalues are the biases
\(\E_{\sigma\in S}(-1)^{\alpha\cdot\sigma}\).

\begin{lemma}[Randomness-efficient amplification]\label{earlier:r:lem:patch-r}
Let \(V_0\) be a PCP verifier that uses \(r_0\) random bits on an
instance \(x\), and let \(0<s\le1\) and \(0<\sigma_0\le1\). Let
\(S\subseteq\{0,1\}^{r_0}\) be \(\varepsilon\)-biased with
\(\varepsilon^2\le s\sigma_0\). Let \(V\) accept \(x\) with proof \(\pi\)
and random string \(R\in\{0,1\}^{r_0}\) exactly when \(V_0\) accepts
\(x\) with \(\pi\) and \(R\oplus\sigma\) for every \(\sigma\in S\).
\begin{enumerate}[label=(\roman*)]
\item \(V\) uses \(r_0\) random bits and \(\abs S\) times as many queries
as \(V_0\); if \(V_0\) is non-adaptive, so is \(V\).
\item A proof that \(V_0\) accepts with probability one is accepted by
\(V\) with probability one.
\item If \(V_0\) rejects \(\pi\) with probability at least \(\sigma_0\),
then \(V\) accepts \(\pi\) with probability at most \(s\).
\end{enumerate}
\end{lemma}
\begin{proof}
Parts (i) and (ii) are immediate, because \(R\oplus\sigma\) is again a
random string of \(V_0\). For (iii), let \(\mathcal B\) be the set of
random strings on which \(V_0\) rejects \(\pi\); its density is
\(\mu\ge\sigma_0\). The verifier \(V\) accepts exactly on the set of
\cref{earlier:r:lem:small-bias}(ii), whose density is at most
\(\varepsilon^2/\mu\le s\sigma_0/\sigma_0=s\).
\end{proof}

The proximity sub-verifiers of~\cite{BCGT13} test whether a function on
an affine subspace of the field is close to the evaluation table of a
polynomial of bounded degree, a Reed--Solomon codeword. Their basic
case is the verifier \(V_{\rm RS,=}\) with its prover \(P_{\rm RS,=}\)
\cite[Algorithms~16 and~7]{BCGT13}, which recurse on subspaces of
about half the dimension. On a subspace of dimension \(\kappa\), the
number of random bits of \(V_{\rm RS,=}\) and the number of field
elements in the proof of \(P_{\rm RS,=}\) are given by the
recursions~(27) and~(28) of~\cite[Section~C.8]{BCGT13}. These
recursions depend on the parameter tuple
\((\eta,\kappa_0,\gamma_{\rm rep},\mu_{\rm rep})\)
of~\cite[Theorem~11.1]{BCGT13}. We use the tuple
\((\eta,\kappa_0,\gamma_{\rm rep},\mu_{\rm rep})=(2,4,0,1)\), which
satisfies the conditions listed in that theorem (Equation~(6)
of~\cite[Theorem~11.3]{BCGT13}).
At this tuple the two recursions read as follows.

\begin{lemma}[One Reed--Solomon test at the tuple \((2,4,0,1)\)]\label{earlier:r:lem:rs-recursion}
Define \(\mathrm{rand}_=(\kappa)=\mathrm{len}_=(\kappa)=0\) for
\(\kappa\le4\) and, for \(\kappa\ge5\),
\[
\begin{aligned}
 \mathrm{rand}_=(\kappa)&=1+\max\bigl\{\ceil{\kappa/2}
 +\mathrm{rand}_=(\floor{\kappa/2}+2),\
 \floor{\kappa/2}+1+\mathrm{rand}_=(\ceil{\kappa/2})\bigr\},\\
 \mathrm{len}_=(\kappa)&=2^{\ceil{\kappa/2}}\cdot2^{\floor{\kappa/2}+2}
 +2^{\ceil{\kappa/2}}\,\mathrm{len}_=(\floor{\kappa/2}+2)
 +2^{\floor{\kappa/2}+1}\,\mathrm{len}_=(\ceil{\kappa/2}).
\end{aligned}
\]
Then for every integer \(\kappa\ge1\),
\[
 \mathrm{rand}_=(\kappa)\le\kappa+3\lgx(2\kappa),\qquad
 \mathrm{len}_=(\kappa)\le2^\kappa\kappa^3.
\]
\end{lemma}
\begin{proof}
For \(\kappa\le4\) both functions vanish, so both bounds hold. We first
show, by strong induction on \(\kappa\ge5\), that for every integer
\(j\) with \(\kappa-4<2^j\),
\begin{equation}\label{earlier:r:eq:rs-induction}
 \mathrm{rand}_=(\kappa)+3\le\kappa+3j,\qquad
 \mathrm{len}_=(\kappa)+2^\kappa\le2^\kappa6^j.
\end{equation}
Here \(j\ge1\), because \(\kappa-4\ge1\). The two arguments
\(a=\floor{\kappa/2}+2\) and \(b=\ceil{\kappa/2}\) of the recursion are
smaller than \(\kappa\). Moreover \(a-4=\floor{(\kappa-4)/2}\) and
\(b-4<a-4\), so \(a-4\) and \(b-4\) are less than \(2^{j-1}\). Hence each
\(c\in\{a,b\}\) satisfies \eqref{earlier:r:eq:rs-induction} with \(j-1\) in place
of \(j\) when \(c\ge5\), by induction. When \(c\le4\), it has
\(\mathrm{rand}_=(c)=\mathrm{len}_=(c)=0\).

\emph{Randomness.} Since \(\ceil{\kappa/2}+a=\kappa+2\) and
\(\floor{\kappa/2}+1+b=\kappa+1\), a branch of the maximum whose
argument \(c\) is at least \(5\) is at most \(\kappa+2+3(j-1)-3\). A
branch with \(c\le4\) is at most \(\floor{\kappa/2}+1\le\kappa-1\). Both
are at most \(\kappa+3j-4\), so
\(\mathrm{rand}_=(\kappa)\le\kappa+3j-3\).

\emph{Length.} By the cases above,
\(\mathrm{len}_=(c)+2^c\le2^c6^{j-1}\) for \(c\in\{a,b\}\). Multiply
the case \(c=a\) by \(2^{\ceil{\kappa/2}}\) and the case \(c=b\) by
\(2^{\floor{\kappa/2}+1}\). Since
\(2^{\ceil{\kappa/2}}\cdot2^a=4\cdot2^\kappa\) and
\(2^{\floor{\kappa/2}+1}\cdot2^b=2\cdot2^\kappa\), adding the two
inequalities gives
\[
 \mathrm{len}_=(\kappa)+2\cdot2^\kappa
 \le6\cdot2^\kappa6^{j-1}=2^\kappa6^j,
\]
which implies \eqref{earlier:r:eq:rs-induction}.

Now let \(\kappa\ge5\), let \(m=\floor{\lgx(\kappa-4)}\), so that
\(2^m\le\kappa-4<2^{m+1}\), and apply \eqref{earlier:r:eq:rs-induction} with
\(j=m+1\). Since \(2^{m+1}\le2(\kappa-4)<2\kappa\), we have
\(m+1\le\lgx(2\kappa)\), and therefore
\(\mathrm{rand}_=(\kappa)\le\kappa+3(m+1)-3\le\kappa+3\lgx(2\kappa)\).
For the length, \(\mathrm{len}_=(\kappa)\le2^\kappa6^{m+1}\), and
\(6^{m+1}\le(2^m+4)^3\le\kappa^3\). The first of these inequalities is
checked directly for \(m\le6\), with equality at \(m=5\). For \(m\ge7\),
\(6^{m+1}\le8^m=(2^m)^3\), since it holds at \(m=7\) and, as \(m\) increases by
one, the two sides grow by the factors \(6\) and \(8\).
\end{proof}

\begin{remark}[The general bound in {\cite[Lemma~C.1]{BCGT13}}]\label{earlier:r:rem:bcgt-c1}
We use \cref{earlier:r:lem:rs-recursion} instead of~\cite[Lemma~C.1]{BCGT13}; at the tuple \((2,4,0,1)\) it gives the same bounds, proved directly from recursions~(27)--(28).
\end{remark}

\subsection{Random-access machines}

The verifier is first simulated by a random-access machine, because
memory accesses of such a machine can be checked by sorting
(\cref{earlier:r:thm:trace}).

\begin{definition}[Random-access machines]\label{earlier:r:def:ram}
A \emph{random-access machine} \(M\) has a word width \(w\ge1\), a
number \(K\) of registers, and a code \((I_0,\ldots,I_{L-1})\) with
\(1\le L\le2^w\). A \emph{word} is an integer in \([0,2^w)\), written
with \(w\) bits. A \emph{configuration} consists of a program counter
\(\mathit{pc}\), which is a word, and registers \(r_0,\ldots,r_{K-1}\),
each a word or the end-of-tape symbol \(\#\). The memory has \(2^w\)
cells, each holding a word. Two read-only tapes, tape~A for the input
and tape~B for the witness, are read in consecutive words of \(w\)
bits; bit \(j\) of a word is the \(j\)-th bit of its block, least
significant first. An operand is a register, \(\mathit{pc}\) or a constant word; a
register holding \(\#\) reads as \(0\). The instructions are:
\begin{itemize}
\item \(\mathtt{readA}\,i\) and \(\mathtt{readB}\,i\): \(r_i\) receives a
word or \(\#\);
\item \(\mathtt{load}\,i\,j\): \(r_j\leftarrow\mathrm{mem}[r_i]\), and
\(\mathtt{store}\,j\,x\): \(\mathrm{mem}[r_j]\leftarrow x\);
\item \(\mathtt{iseof}\,i\,j\): \(r_i\leftarrow1\) if \(r_j=\#\), and
\(r_i\leftarrow0\) otherwise;
\item \(\mathtt{arith}\,o\,i\,x\,y\): \(r_i\leftarrow o(x,y)\) for
\(o\in\{\mathrm{add},\mathrm{sub},\mathrm{and},\mathrm{cmp}\}\), where
addition and subtraction are modulo \(2^w\), \(\mathrm{and}\) is
bitwise, and \(\mathrm{cmp}(x,y)\) is \(1\), \(0\) or \(2^w-1\)
according as \(x<y\), \(x=y\) or \(x>y\);
\item \(\mathtt{nop}\);
\item \(\mathtt{je}\,i\,c\), \(\mathtt{jne}\,i\,c\) and
\(\mathtt{jmp}\,c\): \(\mathit{pc}\leftarrow c\bmod2^w\) if \(r_i=0\),
if \(r_i\ne0\), or always;
\item \(\mathtt{out}\,\mathrm{accept}\) and
\(\mathtt{out}\,\mathrm{reject}\): halt.
\end{itemize}
Every instruction except a taken jump and \(\mathtt{out}\) also sets
\(\mathit{pc}\leftarrow\mathit{pc}+1\bmod2^w\). A configuration whose
\(\mathit{pc}\) is at least \(L\) or points to \(\mathtt{out}\) has no
successor. The initial configuration has \(\mathit{pc}=0\) and all
registers \(0\); a configuration is accepting if
\(I_{\mathit{pc}}=\mathtt{out}\,\mathrm{accept}\).

Let \(T\ge1\). An \emph{accepting run of length \(T\)} on binary
strings \((\bar x,y)\) consists of configurations
\(S^0,\ldots,S^{T-1}\) and memories
\(\mathrm{mem}_0,\ldots,\mathrm{mem}_{T-1}\) such that \(S^0\) is
initial, \(\mathrm{mem}_0=0\), \(S^{T-1}\) is accepting, and
\((S^{i+1},\mathrm{mem}_{i+1})\) results from \((S^i,\mathrm{mem}_i)\)
by executing \(I_{\mathit{pc}}\) for every \(i<T-1\); moreover the
words other than \(\#\) that the run reads from tape~A, concatenated in
time order, form a prefix of \(\bar x\), and those read from tape~B form
a prefix of \(y\). The machine \emph{accepts \(\bar x\) in time \(T\)}
if it has such a run for some \(y\) with \(\abs y\le T\).
\end{definition}

A read may return \(\#\) at any step; the prefix conditions make every
word that is returned the next word of its tape. This frees the trace
from tracking where a tape ends. The machines \(M_n\) never rely on it:
a read of \(\#\) sends them into a loop that never halts
(\cref{earlier:r:lem:tm-ram}).

\begin{lemma}[Multitape machines as random-access machines]\label{earlier:r:lem:tm-ram}
Let \(V(x,y)\) be a deterministic multitape machine that halts within
\(T(n)\ge n\) steps on every \(x\in\{0,1\}^n\) and every \(y\). There are
constants \(c_V,K,L\ge1\) with the following properties. For each
\(n\ge2\) let the word width
\(w_n=2^{\ceil{\lgx(\ceil{\lgx T(n)}+c_V)}}\) be the least power of two
with \(w_n\ge\ceil{\lgx T(n)}+c_V\), for \(x\in\{0,1\}^n\) let
\(\bar x=x10^j\) be its padding to length
\(2^{\ceil{\lgx\max(n+1,w_n)}}\), a multiple of \(w_n\), and let \(T_n\)
be the least power of two with \(T_n\ge c_V(T(n)+w_n)\). There is a
random-access machine \(M_n\) with word width \(w_n\), \(K\) registers
and a code of length \(L\) such that, for every \(x\in\{0,1\}^n\),
\(V(x,y)\) accepts for some \(y\in\{0,1\}^{T(n)}\) if and only if
\(M_n\) accepts \(\bar x\) in time \(T_n\). Moreover
\(\abs{\bar x}\le T_n<2^{w_n}\), and the code of \(M_n\) is a fixed
code, depending only on \(V\), into which the integers \(n\),
\(T(n)\) and \(T_n-8\) are written.
\end{lemma}

The proof of \cref{earlier:r:lem:tm-ram} is a direct simulation followed by an
idle phase that pads the run to exactly \(T_n\) configurations. It is
routine, and we give it at the end of this appendix.

\subsection{Affine constraint instances}

The target of the reduction is the following algebraic problem. Let
\(\mathrm{irr}\in\mathbb F_2[X]\) be irreducible of degree \(f\), and
let \(F=\mathbb F_2[X]/(\mathrm{irr})\), the field with \(2^f\)
elements. An affine subspace of \(F\) over \(\mathbb F_2\) is given by a
basis \(b_1,\ldots,b_m\) and an offset \(o\); its \(i\)-th element,
\(0\le i<2^m\), is \(o+\sum_ji_jb_j\), where \(i_1,i_2,\ldots\) are the
binary digits of \(i\), least significant first.

\begin{definition}[Affine constraint instances]\label{earlier:r:def:acsp}
An \emph{affine constraint instance} consists of an irreducible
\(\mathrm{irr}\) of degree \(f\), an affine subspace \(H\subseteq F\) of
dimension \(m_H\), neighbors \(N_\nu(X)=a_\nu X+b_\nu\) with
\(a_\nu\ne0\), for \(1\le\nu\le c_N\) and some \(c_N\ge1\), a constraint
polynomial \(P(X,Y_1,\ldots,Y_{c_N})\) over \(F\) given by an arithmetic
circuit, and an affine subspace \(I\subseteq H\) of dimension \(m\).
Give \(X\) weight \(1\) and each \(Y_\nu\) weight \(2^{m_H}-1\). The
instance has \emph{degree bound} \(d_B\) if every monomial of \(P\) has
weight at most \(d_B\); we require \(4d_B\le2^f\). The instance
\emph{accepts} a binary string \(\bar x\) of length \(2^m\) if some
\(A\in F[X]\) of degree less than \(2^{m_H}\) satisfies
\[
 P\bigl(\alpha,A(N_1(\alpha)),\ldots,A(N_{c_N}(\alpha))\bigr)=0
 \qquad\text{for every }\alpha\in H,
\]
and \(A(\alpha_i)=\bar x_i\) for every \(i<2^m\), where \(\alpha_i\) is
the \(i\)-th element of \(I\) and the bits \(0,1\) are read in \(F\).
\end{definition}

If \(\deg A<2^{m_H}\), then
\(B_A(X)=P(X,A(N_1(X)),\ldots,A(N_{c_N}(X)))\) has degree at most
\(d_B\): substituting \(A\circ N_\nu\), of degree at most
\(2^{m_H}-1\), raises no monomial above its weight. This is the succinct
algebraic constraint satisfaction problem
of~\cite[Definition~7.1]{BCGT13} at a single instance, with affine
neighbors. The weighted bound \(4d_B\le2^f\) implies its degree
condition~(3) and, as just shown, bounds every composition \(B_A\); we
use this stronger sufficient condition.

\subsection{The run as an event trace}\label{earlier:r:app:event-trace}

\emph{Idea.} A run is checked locally in two orders. In time order,
each step relates two consecutive configurations, which is a local
check, except that a \(\mathtt{load}\) must return the value last stored
at its address, possibly many steps earlier. In address order, with
ties broken by time, every \(\mathtt{load}\) directly follows the last
access to its cell, so the memory check is local again. The trace
therefore lists each memory access, and each word read from the input
tape, as a \emph{record}: once in time order, at the step that makes
it, and once in sorted order. The two lists must be equal as multisets;
in the affine instance a routing network enforces this
(\cref{earlier:r:lem:routing}). Checking random-access memory by sorting is a
standard technique; the argument below is self-contained.

Fix a machine \(M\) with word width \(w=2^\omega\), \(K\) registers and
code length \(L\); an integer \(\tau\ge\max(1,\omega)\) with
\(T=2^\tau<2^w\) (for \(M_n\), \(T=T_n\), not \(T(n)\)); and a string \(\bar x\) whose length is a multiple of
\(w\) and at most \(T\). Put \(\beta=\tau-\omega\),
\(N_A=\abs{\bar x}/w\), \(N=2T\) and \(\kappa_{\rm key}=w+\tau+2\). Every slot
below holds one bit; a \(w\)-bit word \(v\) occupies \(w\) slots
\(v[0],\ldots,v[w-1]\), least significant first.

\begin{definition}[Event trace]\label{earlier:r:def:event-trace}
A \emph{record} consists of an activity bit \(\mathit{act}\), a region
bit \(\mathit{rgn}\) (\(0\) for memory, \(1\) for tape~A), an address or
tape index \(\mathit{ad}<2^w\), a phase bit \(\mathit{ph}\), a time stamp
\(\mathit{ts}<2^\tau\), a value \(\mathit{va}<2^w\), and a store bit
\(\mathit{st}\). Its \emph{key} is
\(\mathit{ts}+2^\tau\mathit{ph}+2^{\tau+1}\mathit{ad}
+2^{\tau+1+w}\mathit{rgn}<2^{\kappa_{\rm key}}\). The active records are memory
accesses \((1,0,a,1,i,v,\mathit{st})\) at address \(a\) and step \(i\)
with value \(v\), tape-A reads \((1,1,c,1,i,v,0)\) of the \(c\)-th word
read, at step \(i\), and input words \((1,1,c,0,0,v,0)\), word \(c\) of
\(\bar x\). The \emph{sentinel} is \((1,0,0,0,0,0,0)\).

The \emph{event trace} \(\mathrm{ET}(M,\tau,\bar x)\) is the following
system of bit slots and equations over \(\mathbb F_2\).
\begin{itemize}
\item \(T\) \emph{time rows} \(\rho_0,\ldots,\rho_{T-1}\). Row \(\rho_i\)
holds the current and the next configuration (\(\mathit{pc}\), and for
each register \(r\) a flag \(\mathit{fl}_r\), set for \(\#\), and a word
\(\mathit{rg}_r\); next values are primed), a one-hot instruction vector
\(e_0,\ldots,e_{L-1}\) with its prefix-or \(o_0,\ldots,o_{L-1}\), operand
words \(X,Y\) and a flag \(X_F\), the results \(\mathit{AD},\mathit{SB},
\mathit{AN}\) of addition, subtraction and conjunction with carries
\(\mathit{ca}_j\), borrows \(\mathit{bo}_j\), a prefix-or
\(\mathit{oz}_j\) and a bit \(\mathit{gt}\), a result word \(Z\) with flag
\(Z_F\), a tape word \(R_V\) with flag \(R_F\), a memory word \(M_V\),
a record \(\mathrm{REC}\) whose time stamp \(\mathit{ts}\) is the row's
time stamp, its successor \(\mathit{ts}'\), the counts \(c_A\) and
\(c_B\) of words read so far from tapes A and~B (\(\tau\) and
\(\beta+1\) bits) with successors \(c_A',c_B'\), increments \(a_A,a_B\),
and the carries \(\mathit{tc}_j,\mathit{ac}_j,\mathit{bc}_j\) of the
three counters.
\item \(N\) \emph{sorted vertices} \(\sigma_0,\ldots,\sigma_{N-1}\), each
with a record \(\mathrm{REC}\), a copy \(\mathrm{PREV}\) of the first
six fields of the previous record, written \(\mathit{pact}\),
\(\mathit{prgn}\), \(\mathit{pad}\), \(\mathit{pph}\), \(\mathit{pts}\)
and \(\mathit{pva}\), and auxiliary bits \(\mathit{both}\),
\(\mathit{ob}_l\), \(\mathit{so}_l\), \(u\), \(g\), \(\mathit{EV}\),
\(g_{A1}\) and \(g_A\).
\item \(N_A\) \emph{input records} \(\iota_0,\ldots,\iota_{N_A-1}\) and
\(N-T-N_A\) \emph{dummy records}, each a record only.
\end{itemize}
The \emph{source list} is \((\mathrm{REC}(\rho_0),\ldots,
\mathrm{REC}(\rho_{T-1}),\iota_0,\ldots,\iota_{N_A-1})\) followed by the
dummies; it has length \(N\). The equations are listed below:
(T1)--(T10) at every time row, (S1)--(S5) at every sorted vertex, the
boundary and copy equations (B1)--(B5), and
\begin{enumerate}[label=(\Alph*),start=7]
\item every slot is a bit, and
\setcounter{enumi}{15}
\item the source list and \((\mathrm{REC}(\sigma_0),\ldots,
\mathrm{REC}(\sigma_{N-1}))\) are equal as multisets.
\end{enumerate}
The system is \emph{satisfiable} if some assignment of bits to all
slots satisfies every equation.
\end{definition}

\emph{Code tables.} Let \(e(Q)=\sum_{p\in Q}e_p\) for a set
\(Q\subseteq[0,L)\) of code positions, a linear form in the slots. The
sets used are \(\mathrm{Ld}\), \(\mathrm{St}\), \(\mathrm{RdA}\),
\(\mathrm{RdB}\), \(\mathrm{Eof}\), \(\mathrm{Jmp}\), \(\mathrm{Je}\),
\(\mathrm{Jne}\), \(\mathrm{Out}\) and \(\mathrm{Ar}_o\) (the positions of
\(\mathtt{load}\), \(\mathtt{store}\), \(\mathtt{readA}\),
\(\mathtt{readB}\), \(\mathtt{iseof}\), the three jumps, \(\mathtt{out}\)
and \(\mathtt{arith}\,o\)), \(\mathrm{Rd}=\mathrm{RdA}\cup\mathrm{RdB}\),
\(\mathrm{Acc}\) (the positions of \(\mathtt{out}\,\mathrm{accept}\)),
\(\mathrm{Seq}\) (the positions of all instructions other than jumps and
\(\mathtt{out}\)), and \(\mathrm{Wr}_r\) (the positions of the
instructions that write register \(r\)). For \(p<L\) let
\(p^+=p+1\bmod2^w\), let \(c_p=c\bmod2^w\) for a jump to \(c\), and let
\(s^X_p[j]\), \(s^Y_p[j]\) be bit \(j\) of the operands that \(I_p\)
places on the buses: \((x,y)\) for \(\mathtt{arith}\,o\,i\,x\,y\);
\((r_i,0)\) for \(\mathtt{je}\) and \(\mathtt{jne}\) on \(r_i\); \(r_i\)
on \(X\) for \(\mathtt{load}\,i\,j\); \((r_j,x)\) for
\(\mathtt{store}\,j\,x\); and \(0\) otherwise. Each is a current slot or
a constant bit. All sums are in \(\mathbb F_2\), and \(j\) ranges over
the bit positions of the word concerned.

\emph{Time rows.}
\begin{enumerate}[label=(T\arabic*)]
\item \(\mathit{fl}_r\,\mathit{rg}_r[j]=0\).
\item \(o_0=e_0\); \(o_p=o_{p-1}+e_p\) and \(e_po_{p-1}=0\) for
\(1\le p<L\); \(o_{L-1}=1\); \(\mathit{pc}[j]=\sum_p\mathrm{bit}_j(p)e_p\).
\item \(X[j]=\sum_pe_ps^X_p[j]\), \(Y[j]=\sum_pe_ps^Y_p[j]\), and
\(X_F=\sum e_p\mathit{fl}_{j'}\) over the positions \(p\) with
\(I_p=\mathtt{iseof}\,i\,j'\).
\item With \(\mathit{ca}_0=\mathit{bo}_0=0\):
\(\mathit{AD}[j]=X[j]+Y[j]+\mathit{ca}_j\) and
\(\mathit{ca}_{j+1}=X[j]Y[j]+X[j]\mathit{ca}_j+Y[j]\mathit{ca}_j\)
(\(j+1<w\)); \(\mathit{SB}[j]=X[j]+Y[j]+\mathit{bo}_j\) and
\(\mathit{bo}_{j+1}=(1+X[j])(Y[j]+\mathit{bo}_j)+Y[j]\mathit{bo}_j\);
\(\mathit{AN}[j]=X[j]Y[j]\); \(\mathit{oz}_0=\mathit{SB}[0]\) and
\(\mathit{oz}_j=\mathit{oz}_{j-1}+\mathit{SB}[j]+\mathit{oz}_{j-1}
\mathit{SB}[j]\); \(\mathit{gt}=(1+\mathit{bo}_w)\mathit{oz}_{w-1}\).
\item \(R_FR_V[j]=0\) and \(e(\mathrm{St})(M_V[j]+Y[j])=0\).
\item \(Z[j]=e(\mathrm{Ar}_{\rm add})\mathit{AD}[j]
+e(\mathrm{Ar}_{\rm sub})\mathit{SB}[j]+e(\mathrm{Ar}_{\rm and})
\mathit{AN}[j]+e(\mathrm{Ar}_{\rm cmp})\mathit{CMP}_j+e(\mathrm{Ld})
M_V[j]+e(\mathrm{Rd})R_V[j]+[j=0]\,e(\mathrm{Eof})X_F\), where
\(\mathit{CMP}_0=\mathit{bo}_w+\mathit{gt}\) and
\(\mathit{CMP}_j=\mathit{gt}\) for \(j\ge1\); and
\(Z_F=e(\mathrm{Rd})R_F\).
\item \(\mathit{rg}_r'[j]=e(\mathrm{Wr}_r)Z[j]+(1+e(\mathrm{Wr}_r))
\mathit{rg}_r[j]\) and \(\mathit{fl}_r'=e(\mathrm{Wr}_r)Z_F
+(1+e(\mathrm{Wr}_r))\mathit{fl}_r\).
\item With \(z=\mathit{oz}_{w-1}\),
\begin{multline*}
 \mathit{pc}'[j]=\sum_{\mathrm{Seq}}e_p\mathrm{bit}_j(p^+)
 +\sum_{\mathrm{Out}}e_p\mathrm{bit}_j(p)
 +\sum_{\mathrm{Jmp}}e_p\mathrm{bit}_j(c_p)\\
 +\sum_{\mathrm{Je}}\bigl(e_pz\,\mathrm{bit}_j(p^+)
 +(e_p+e_pz)\mathrm{bit}_j(c_p)\bigr)
 +\sum_{\mathrm{Jne}}\bigl(e_pz\,\mathrm{bit}_j(c_p)
 +(e_p+e_pz)\mathrm{bit}_j(p^+)\bigr).
\end{multline*}
\item \(\mathit{act}=e(\mathrm{Ld})+e(\mathrm{St})+a_A\),
\(\mathit{rgn}=a_A\),
\(\mathit{ad}[j]=(e(\mathrm{Ld})+e(\mathrm{St}))X[j]+a_Ac_A[j]\)
(with \(c_A[j]=0\) for \(j\ge\tau\)), \(\mathit{ph}=1\),
\(\mathit{va}[j]=(e(\mathrm{Ld})+e(\mathrm{St}))M_V[j]+a_AR_V[j]\),
\(\mathit{st}=e(\mathrm{St})\).
\item With \(\mathit{tc}_0=1\): \(\mathit{ts}'[j]=\mathit{ts}[j]
+\mathit{tc}_j\) and \(\mathit{tc}_{j+1}=\mathit{ts}[j]\mathit{tc}_j\).
\(a_A=e(\mathrm{RdA})(1+R_F)\) and, with \(\mathit{ac}_0=a_A\),
\(c_A'[j]=c_A[j]+\mathit{ac}_j\) and
\(\mathit{ac}_{j+1}=c_A[j]\mathit{ac}_j\). Likewise
\(a_B=e(\mathrm{RdB})(1+R_F)\) and \(c_B'=c_B+a_B\) with carries
\(\mathit{bc}_j\). Finally \(a_Bc_B[\beta]=0\).
\end{enumerate}
The last carries of \(\mathit{ts}'\), \(c_A'\) and \(c_B'\) are not
kept, so these counters work modulo \(2^\tau\), \(2^\tau\) and
\(2^{\beta+1}\). In words: (T1) a register holding \(\#\) has word
\(0\); (T2) \(\mathit{pc}\) is decoded one-hot; (T3) the operand buses;
(T4) sum, difference, conjunction and comparison; (T5) the read and
store channels; (T6) the result; (T7) the register write; (T8) the next
\(\mathit{pc}\); (T9) the row's record; (T10) the time stamp and the read
counters.

\emph{Sorted vertices.} Let \(K_l\) and \(P_l\), \(l<\kappa_{\rm key}\), be the key
bits of \(\mathrm{REC}\) and of \(\mathrm{PREV}\), least significant
first.
\begin{enumerate}[label=(S\arabic*)]
\item \((1+\mathit{pact})\mathit{act}=0\) and
\(\mathit{both}=\mathit{act}\cdot\mathit{pact}\).
\item With \(\mathit{ob}_0=0\):
\(\mathit{ob}_{l+1}=(1+P_l)(K_l+\mathit{ob}_l)+K_l\mathit{ob}_l\) for
\(l<\kappa_{\rm key}\), and \(\mathit{both}(1+\mathit{ob}_{\kappa_{\rm key}})=0\).
\item With \(d_0=\mathit{prgn}+\mathit{rgn}\) and
\(d_{1+j}=\mathit{pad}[j]+\mathit{ad}[j]\): \(\mathit{so}_0=d_0\) and
\(\mathit{so}_l=\mathit{so}_{l-1}+d_l+\mathit{so}_{l-1}d_l\) for
\(1\le l\le w\).
\item \(u=\mathit{act}(1+\mathit{rgn})\), \(g=u(1+\mathit{st})\),
\(\mathit{EV}[j]=(1+\mathit{so}_w)\mathit{pva}[j]\) and
\(g(\mathit{va}[j]+\mathit{EV}[j])=0\).
\item \(g_{A1}=\mathit{act}\cdot\mathit{rgn}\),
\(g_A=g_{A1}\mathit{ph}\), \(g_A\mathit{so}_w=0\),
\(g_A\mathit{pph}=0\) and \(g_A(\mathit{va}[j]+\mathit{pva}[j])=0\).
\end{enumerate}
In words: (S1) active records come first; (S2) keys increase; (S3) the
record names the same region and address as its predecessor exactly
when \(\mathit{so}_w=0\); (S4) the memory rule; (S5) the input rule.

\emph{Boundaries and copies.}
\begin{enumerate}[label=(B\arabic*)]
\item At \(\rho_0\): \(\mathit{pc}=0\), \(\mathit{fl}_r=0\) and
\(\mathit{rg}_r=0\) for all \(r\), \(\mathit{ts}=0\), \(c_A=0\),
\(c_B=0\).
\item For \(i<T-1\): \(e(\mathrm{Out})=0\) at \(\rho_i\), and
\(\rho_{i+1}\) has current configuration, \(\mathit{ts}\), \(c_A\) and
\(c_B\) equal to the next configuration, \(\mathit{ts}'\), \(c_A'\)
and \(c_B'\) of \(\rho_i\).
\item At \(\rho_{T-1}\): \(e(\mathrm{Acc})=1\).
\item \(\mathrm{PREV}(\sigma_0)\) is the sentinel's first six fields,
and \(\mathrm{PREV}(\sigma_{j+1})\) is the first six fields of
\(\mathrm{REC}(\sigma_j)\).
\item \(\iota_c=(1,1,c,0,0,v_c,0)\) with \(v_c\) word \(c\) of
\(\bar x\); every dummy record is zero.
\end{enumerate}

Every equation has degree at most two in the slots; the auxiliary slots
\(Z_F\), \(a_A\), \(a_B\), \(u\), \(g\), \(g_{A1}\), \(g_A\),
\(\mathit{EV}\), \(\mathit{both}\), \(\mathit{oz}_j\) and \(\mathit{gt}\)
exist to keep it so.

\begin{lemma}[Gadgets]\label{earlier:r:lem:trace-gadgets}
Assume (G). For all words \(X,Y\), the equations (T4) have exactly one
solution in the auxiliary slots: the carries, borrows and prefix-ors of
the school methods. In it \(\mathit{AD}=X+Y\bmod2^w\),
\(\mathit{SB}=X-Y\bmod2^w\), \(\mathit{bo}_w=[X<Y]\),
\(\mathit{AN}=X\wedge Y\), \(\mathit{oz}_{w-1}=[X\ne Y]\),
\(\mathit{gt}=[X>Y]\) and \(\mathit{CMP}=\mathrm{cmp}(X,Y)\). Likewise
(T10) forces \(\mathit{ts}'=\mathit{ts}+1\bmod2^\tau\),
\(c_A'=c_A+a_A\bmod2^\tau\) and \(c_B'=c_B+a_B\bmod2^{\beta+1}\), with
unique carries.
\end{lemma}
\begin{proof}
Each auxiliary slot is the left side of one equation whose right side
uses inputs and earlier slots only, so the solution exists and is
unique. By induction on \(j\),
\(\mathit{ca}_j=\floor{(X\bmod2^j+Y\bmod2^j)/2^j}\), because over
\(\mathbb F_2\) the carry equation is the majority of three bits;
likewise \(\mathit{bo}_j=[X\bmod2^j<Y\bmod2^j]\), the majority of
\(1+X[j]\), \(Y[j]\) and \(\mathit{bo}_j\), and
\(\mathit{oz}_j=[\mathit{SB}\bmod2^{j+1}\ne0]\). Since \(\mathit{SB}=0\)
exactly when \(X=Y\), \(\mathit{gt}=[X\ge Y\text{ and }X\ne Y]\). The
counters are the case \(Y\in\{0,1\}\) of the adder.
\end{proof}

\begin{lemma}[One step]\label{earlier:r:lem:trace-step}
Assume (G).
\begin{enumerate}[label=(\roman*)]
\item (T1) holds for a configuration block exactly when it encodes a
configuration: a register is \(\#\) if its flag is set, and its word is
then \(0\).
\item (T2) holds exactly when \(\mathit{pc}<L\), \(e\) is the indicator
vector of \(\mathit{pc}\), and \(o_p=[\mathit{pc}\le p]\).
\item Let a time row satisfy (T1)--(T10), let \(S\) and \(S'\) be its
current and next configurations, and let \(I=I_{\mathit{pc}}\). If \(I\)
is not \(\mathtt{out}\), then \(S'\) results from \(S\) by executing
\(I\) with the row's access: a \(\mathtt{load}\) receives \(M_V\) from
address \(X\), a \(\mathtt{store}\) writes \(M_V=Y\) at address \(X\),
and a read receives \(\#\) if \(R_F=1\) and \(R_V\) otherwise. If \(I\)
is \(\mathtt{out}\), then \(S'=S\). Conversely, every such step, with
the auxiliary slots of \cref{earlier:r:lem:trace-gadgets}, unused channels zero,
and \(c_B[\beta]=0\) whenever \(I\) is a \(\mathtt{readB}\) that receives a
word, satisfies (T1)--(T10).
\end{enumerate}
\end{lemma}
\begin{proof}
(i) restates the encoding. (ii) The equations make \(o_p\) the number
of ones among \(e_0,\ldots,e_p\), keep it at most one, and make it one
at \(p=L-1\); then \(\mathit{pc}=\sum_pp\,e_p\). In characteristic two,
\(\sum_pe_p=1\) alone would allow three ones; the prefix-or chain
excludes this. (iii) By (ii) every \(e(Q)\) equals
\([\mathit{pc}\in Q]\). (T3) puts the operands of \(I\) on the buses,
\cref{earlier:r:lem:trace-gadgets} computes the four operations and the zero test
\(z=[X\ne0]\) used by the conditional jumps, (T6) selects the value
written, (T7) writes it into the register of \(I\) and copies the
others, and (T8) computes the next \(\mathit{pc}\). A read of \(\#\)
writes the flag \(1\) and the word \(0\) by (T5) and (T6). The converse
is the same computation.
\end{proof}

\begin{lemma}[Memory by sorting]\label{earlier:r:lem:trace-memory}
Let a time-ordered list of memory accesses be given, at distinct steps.
The following are equivalent.
\begin{enumerate}[label=(\alph*)]
\item There are memories \(\mathrm{mem}_0=0,\mathrm{mem}_1,\ldots\) such
that each store \((a,v)\) at step \(i\) gives
\(\mathrm{mem}_{i+1}=\mathrm{mem}_i[a\mapsto v]\), every other step
leaves the memory unchanged, and each load \((a,v)\) at step \(i\) has
\(v=\mathrm{mem}_i(a)\).
\item Each load \((a,v)\) at step \(i\) has \(v\) equal to the value of
the latest access to \(a\) before step \(i\), or \(v=0\) if there is
none.
\item In the list of the memory records sorted by key and preceded by
the sentinel, a load record whose predecessor has the same address has
the predecessor's value, and every other load record has value \(0\).
\end{enumerate}
\end{lemma}
\begin{proof}
(a)\(\Leftrightarrow\)(b): by induction on \(i\), \(\mathrm{mem}_i(a)\) is
the value of the latest access to \(a\) before step \(i\), or \(0\).
(b)\(\Leftrightarrow\)(c): the memory record of an access to \(a\) at
step \(i\) has key \(2^{\tau+1}a+2^\tau+i\). In key order the records of
one address are therefore consecutive and in time order. The
predecessor of a load at \((a,i)\) is the latest access to \(a\) before
\(i\) if there is one, and otherwise a record of a smaller address or
the sentinel, whose address \(0\) and value \(0\) give the value \(0\)
in either case.
\end{proof}

\begin{lemma}[Time stamps and counters]\label{earlier:r:lem:trace-counters}
Assume (G), (T1)--(T10) at every time row, (B1) and (B2). Then
\(\mathit{ts}(\rho_i)=i\), and \(c_A(\rho_i)\) and \(c_B(\rho_i)\) are
the numbers of steps before \(i\) that read a word other than \(\#\)
from tape~A and from tape~B. Hence at most \(2^\beta\) words are read
from tape~B, of total length at most \(2^\beta w=T\). Conversely, if at
most \(2^\beta\) such words are read, the honest counters satisfy
\(a_Bc_B[\beta]=0\).
\end{lemma}
\begin{proof}
Induction on \(i\) with \cref{earlier:r:lem:trace-gadgets}. The time stamp does
not wrap, since \(i\le T-1<2^\tau\); nor does \(c_A\), since at most
\(T-1\) steps precede \(\rho_{T-1}\). A read at \(c_B=2^\beta\) is
excluded by \(a_Bc_B[\beta]=0\), so \(c_B\le2^\beta<2^{\beta+1}\).
Conversely the \(r\)-th read happens at \(c_B=r-1<2^\beta\).
\end{proof}

\begin{lemma}[The sorted list]\label{earlier:r:lem:trace-sorted}
Assume (G), (P), (S1)--(S5), (B4), (B5), the time-row equations, and
\(\mathit{ts}(\rho_i)=i\) for all \(i\). Then the active sorted vertices
form a prefix of \(\sigma_0,\ldots,\sigma_{N-1}\), which lists the active
records of the source list in increasing key order. The memory records
satisfy \cref{earlier:r:lem:trace-memory}(c), and every tape-A read record
\((1,1,c,1,i,v,0)\) has \(c<N_A\) and \(v\) equal to word \(c\) of
\(\bar x\). Conversely, let a source list have distinct active keys and
pairwise distinct tape-A read indices \(c\), and let its memory and
tape-A records have these two properties. Then sorting its active
records by key, appending the inactive ones, and filling
\(\mathrm{PREV}\) by (B4) satisfies (S1)--(S5) and (P).
\end{lemma}
\begin{proof}
By (S1) activity does not increase along \(\sigma\), and the sentinel is
active, so the active vertices form a prefix. By (S2) their keys
increase strictly; by (P) they are the active source records in key
order, and their keys are distinct. A memory load has \(u=g=1\), and
(S4) is \cref{earlier:r:lem:trace-memory}(c), because \(\mathit{so}_w=0\) exactly
when the two records name the same region and address. A tape-A read
\((1,1,c,1,i,v,0)\) has \(g_A=1\), and (S5) makes its predecessor a
record in region~\(1\) with index \(c\), phase \(0\) and value \(v\).
Time rows have phase \(1\) by (T9), so the predecessor is \(\iota_c\),
which exists only for \(c<N_A\) and has the value prescribed by (B5).
For the converse, each predecessor in the sorted arrangement is as
described, and the same computation verifies every equation. Distinct
read indices are needed here: a second read of index \(c\) would have
a read record of phase~\(1\) as its predecessor and violate (S5).
\end{proof}

\begin{theorem}[Event traces]\label{earlier:r:thm:trace}
\(\mathrm{ET}(M,\tau,\bar x)\) is satisfiable if and only if \(M\)
accepts \(\bar x\) in time \(T=2^\tau\).
\end{theorem}
\begin{proof}
(\(\Rightarrow\)) Let \(S^i\) be the current configuration of \(\rho_i\)
(\cref{earlier:r:lem:trace-step}(i)). By (B2) the next configuration of \(\rho_i\)
is \(S^{i+1}\); by (B1) \(S^0\) is initial; by (B2) and (B3) no
configuration before the last executes \(\mathtt{out}\), and
\(S^{T-1}\) is accepting. \Cref{earlier:r:lem:trace-step}(iii) gives every step
relative to its access, and \cref{earlier:r:lem:trace-counters} gives the time
stamps, so the access records have distinct keys. By
\cref{earlier:r:lem:trace-sorted,earlier:r:lem:trace-memory}, there are memories with
\(\mathrm{mem}_0=0\) under which every step is a step of \(M\). The
words read from tape~A have indices \(c=0,1,2,\ldots\) in time order
(\cref{earlier:r:lem:trace-counters}), and each equals word \(c\) of \(\bar x\);
so they form a prefix of \(\bar x\). Let \(y\) be the concatenation of
the words read from tape~B; then \(\abs y\le T\) by
\cref{earlier:r:lem:trace-counters}, and the run is accepting on \((\bar x,y)\).

(\(\Leftarrow\)) Take an accepting run on \((\bar x,y)\) with
\(\abs y\le T\). Fill \(\rho_i\) from \(S^i\), \(S^{i+1}\) and the
access of step \(i\) (the loaded value \(\mathrm{mem}_i(a)\), the stored
word, or the word read), with the auxiliary slots of
\cref{earlier:r:lem:trace-gadgets}, unused channels zero, the next configuration of \(\rho_{T-1}\) equal
to its current one, and honest time stamps and counters. At most \(\abs{\bar x}/w=N_A\) words
are read from tape~A, and the \(c\)-th is word \(c\) of \(\bar x\); at
most \(\abs y/w\le2^\beta\) words are read from tape~B, so the witness
bound holds (\cref{earlier:r:lem:trace-counters}). Add the input and dummy
records. The active keys are distinct: time-row records differ in their
time stamps, input records in their indices, and the two kinds in their
phase. The tape-A read indices are the counter values
\(0,1,\ldots\). \Cref{earlier:r:lem:trace-sorted} now provides the sorted
vertices, and every equation holds by
\cref{earlier:r:lem:trace-gadgets,earlier:r:lem:trace-step,earlier:r:lem:trace-memory,earlier:r:lem:trace-counters}.
\end{proof}

\begin{corollary}[Event traces of verifiers]\label{earlier:r:cor:trace-verifier}
Let \(V\), \(T(n)\) and \(M_n\) be as in \cref{earlier:r:lem:tm-ram}, let
\(n\ge2\), and let \(T_n=2^{\tau_n}\). For every \(x\in\{0,1\}^n\),
\(V(x,y)\) accepts for some \(y\in\{0,1\}^{T(n)}\) if and only if
\(\mathrm{ET}(M_n,\tau_n,\bar x)\) is satisfiable.
\end{corollary}
\begin{proof}
\Cref{earlier:r:lem:tm-ram} and \cref{earlier:r:thm:trace}. The hypotheses of the latter
hold: \(\tau_n\ge\max(1,\omega)\) and \(\abs{\bar x}\le T_n<2^{w_n}\) by
\cref{earlier:r:lem:tm-ram}, and \(w_n\) divides \(\abs{\bar x}\).
\end{proof}

\subsection{The event trace as an affine constraint instance}\label{earlier:r:app:affine-embedding}

Throughout this subsection and the next, assume in addition that
\(\abs{\bar x}=2^m\) for an integer \(m\) with \(\omega\le m\le\tau\).
This holds for the padded inputs of \cref{earlier:r:lem:tm-ram}; the event-trace
theorem itself does not need it.

\emph{Idea.} Each slot of the event trace becomes a point of an affine
subspace \(H\) of a binary field, and the assignment polynomial \(A\)
holds its value there. Each equation, at each vertex, has its own host point in \(H\); a
\emph{selector} polynomial switches the equation on at its host points
only, and the equation reads the slots it needs at affine shifts
\(\alpha+b\) of the host point \(\alpha\). Three kinds of reads occur: within a
vertex, from a time row to its predecessor, and from a sorted vertex to
its predecessor. The multiset equality (P) is not local. A Beneš
network~\cite{Ben65} enforces it: between the source list and the sorted list there
are \(2k_N-1\) layers of switches, each of which passes a pair of
records straight or swaps them.

\emph{Coordinates.} Let \(k_N=\tau+1\), so that \(N=2^{k_N}\); let
\(\ell_S=\ceil{\lgx(2k_N)}\); let \(F_{\rm tags}=32+3K+2L\) (the number
of field tags the placement below needs),
\(d=\ceil{\lgx F_{\rm tags}}\), \(q=d+\omega\) and \(C=2^q\); and let
\(m_H=k_N+\ell_S+q+4\). Given an irreducible \(\mathrm{irr}\) of degree
\(f\ge m_H\), the residues \(e_b=X^b\), \(b<m_H\), are linearly
independent over \(\mathbb F_2\). Let \(H\) be their span, with offset
\(0\), so that \(\abs H=2^{m_H}=16NC2^{\ell_S}\). A point of \(H\) is
labeled by a lane \(x<N\), a stage \(s<2^{\ell_S}\), a column \(c<C\), a
role \(r<8\) and an area \(a<2\):
\[
 E(x,s,c,r,a)=\sum_{b<k_N}x_be_b+\sum_{b<\ell_S}s_be_{k_N+b}
 +\sum_{b<q}c_be_{k_N+\ell_S+b}+\sum_{b<3}r_be_{k_N+\ell_S+q+b}
 +a\,e_{m_H-1},
\]
where \(x_b\) is bit \(b\) of \(x\), and so on. A column \(c=2^\omega
t+j\) names a field tag \(t<2^d\) and a bit index \(j<w\). Labels add by
the bitwise exclusive-or of their coordinates.

\emph{Placement.} A fixed catalog, which depends only on \(K\) and
\(L\), assigns every slot and every equation its own point.
\begin{itemize}
\item \emph{Area~0, the network.} At stage \(s<2k_N\) and lane \(x\),
role~\(0\) holds a network record of \(C\) bits, one per column: the
seven record fields occupy their own field tags, the value field
\(\mathit{va}\) at tag~\(0\) (the input subspace below relies on this),
at bit index \(j\) for bit \(j\), and all other columns are padding. Stage~\(0\) holds the
source list (\(\mathrm{REC}(\rho_i)\) at lane \(i<T\), \(\iota_c\) at
lane \(T+c\), dummies at the other lanes), and stage \(2k_N-1\) holds
\(\mathrm{REC}(\sigma_j)\) at lane \(j\). Role~\(1\) at a stage
\(s<2k_N-1\) holds that stage's switch at the lane, copied in every
column. The routing equation of a record bit is hosted at that bit's
own point, the column equation of a switch at the switch's point, the
bit check of a record bit at role~\(2\), the bit check of a switch at
role~\(3\), and the pair equation at role~\(4\).
\item \emph{Area~1, the machine.} The other slots of \(\rho_i\) are at
stage~\(0\), lane \(i\), and those of \(\sigma_j\) at stage \(2k_N-1\),
lane \(j\). Each word-sized block of a vertex has its own field tag;
the blocks use at most \(31+2K+2L\) tags.
\begin{itemize}
\item Roles~\(0\) and~\(1\) hold data: current values in role~\(0\);
next values, and \(\mathrm{PREV}\), in role~\(1\), in the column of the
value they are copied to.
\item Role~\(2\) hosts the equations, one field tag per equation
\emph{template} and bit index \(j\) for the instance at \(j\). A template
is one displayed equation of (T1)--(T10) or (S1)--(S5) for a fixed
register \(r\) or code position \(p\); the \(\kappa_{\rm key}\)-bit borrow chain of
(S2) is split into three templates of at most \(w\) instances each,
since \(\kappa_{\rm key}\le2w+1\). There are exactly \(32+3K+2L\) time templates
and \(17\) sorted-vertex templates, so \(F_{\rm tags}\) tags suffice.
The initial carries \(\mathit{ca}_0\), \(\mathit{bo}_0\),
\(\mathit{tc}_0\), \(\mathit{ac}_0\) and \(\mathit{bc}_0\) are
shorthand for their displayed values, requiring no additional slots or
templates. The respective template counts for (T1)--(T10) are
\(K,\,2L+1,\,3,\,8,\,2,\,2,\,2K,\,1,\,6,\,9\), whose sum is
\(32+3K+2L\).
\item Roles~\(3\) and~\(4\) host the bit checks (G) of roles~\(0\)
and~\(1\), in the same column.
\item Role~\(5\) hosts (B1) and the time copies of (B2), role~\(6\)
hosts (B4), and role~\(7\) hosts (B3) and (B5).
\end{itemize}
\end{itemize}

\emph{Neighbors.} Every neighbor is a translation \(X\mapsto X+b\) with
\(b\in H\), so it maps \(H\) onto \(H\). An equation hosted at a point
\(\alpha\) reads its slots at \(\alpha+b\) for the following \(b\).
\begin{itemize}
\item \emph{Within a vertex:} the \(16C\) sums
\(E(0,0,c,r,a)\) over all columns, roles and areas, which change no
lane or stage bit; the exclusive-or of the host's and the target's
column, role and area selects the target.
\item \emph{Previous time row and previous sorted vertex:} for each
\(t<k_N\), the translation by \(\Gamma_t=e_0+\cdots+e_t\) plus a fixed
role and area mask. If \(t\) is the least set bit of the lane \(x\), then
\(x\oplus(2^{t+1}-1)=x-1\), so this translation reads lane \(x-1\) at
the same column. Two such families of \(k_N\) maps are used.
\item \emph{Network edges:} the identity; for each stage \(s<2k_N-1\),
the translation \(\delta_s\) that changes the stage label \(s\) to
\(s+1\) and \(\delta_s+e_{d_s}\), which in addition flips lane bit
\(d_s=\min(s,2k_N-2-s)\); the role masks used by the routing equations;
the pair-partner translations \(E(0,0,0,5,0)+e_b\), \(b<k_N\), by which
the pair equation at role~\(4\) reads the switch at role~\(1\) of lane
\(x\oplus2^b\); the column successors \(\zeta_t=e_{k_N+\ell_S}+\cdots+e_{k_N+\ell_S+t}\),
\(t<q\); and the lane flips \(e_b\), \(b<k_N\) (counted among the neighbors; no equation of this construction reads them).
\end{itemize}
Hence \(c_N\le8k_N+q+16C+3\). The identity is a neighbor, so
\(c_N\ge1\).

\emph{Selectors.} For each coordinate \(a<m_H\) there is a
\emph{linearized} polynomial
\(\lambda_a(X)=\sum_{r<m_H}c_{a,r}X^{2^r}\), of degree at most
\(2^{m_H-1}\), with \(\lambda_a(e_b)=[a=b]\); since \(\lambda_a\) is
\(\mathbb F_2\)-linear, \(\lambda_a(\sum_bu_be_b)=u_a\) for all bits
\(u_b\). Indeed, the coefficients solve an \(m_H\times m_H\) linear
system with matrix \((e_b^{2^r})_{b,r}\). A nonzero vector in its kernel
would give a nonzero linearized polynomial of degree at most
\(2^{m_H-1}\) vanishing on every \(e_b\), hence on all \(2^{m_H}\) points
of \(H\), which is impossible; so the system is invertible. A
\emph{pattern} is an assignment of bits to a set of distinct
coordinates. It is emitted as
\[
 \chi_p(X)=\prod_{a:\,p(a)=1}\lambda_a(X)\prod_{a:\,p(a)=0}
 \bigl(1+\lambda_a(X)\bigr),
\]
which on \(H\) is the indicator of the points whose coordinates agree
with \(p\), and has degree at most \(m_H2^{m_H-1}\). The selector of an
equation is the sum of a list of disjoint patterns: one pattern fixes
the host's stage, column, role and area, and it is merged with the
patterns of a lane condition. The lane conditions used are the time
lanes, the input lanes, the dummy lanes, the first lane, the last time
lane, the other time lanes, and, for each \(t\), the lanes whose least
set bit is \(t\); each is a list of at most \(k_N\) patterns. Two patterns
are merged by uniting their assignments before emission, and a conflict
emits nothing. Every emitted pattern thus fixes each coordinate at most
once.

\emph{Routing equations.} Write \(D_s(x,c)\) for the network record bit
at stage \(s\), lane \(x\), column \(c\), and \(W_s(x,c)\) for the switch
value. At every stage \(s<2k_N-1\), every lane and every column, the
host points of area~\(0\) carry the equations
\begin{gather*}
 D_s(x,c)+D_{s+1}(x,c)+W_s(x,c)\bigl(D_{s+1}(x,c)
 +D_{s+1}(x\oplus2^{d_s},c)\bigr)=0,\\
 W_s(x,c)^2+W_s(x,c)=0,\\
 W_s(x,c)+W_s(x,c+1\bmod C)=0,\qquad
 W_s(x,c)+W_s(x\oplus2^{d_s},c)=0,
\end{gather*}
and every network record bit satisfies \(D_s(x,c)^2+D_s(x,c)=0\). The
column successor \(c\mapsto c+1\bmod C\) is the translation by
\(\zeta_t\) at the columns whose lowest zero bit is \(t\), and by
\(\zeta_{q-1}\) at the column of all ones; each case is selected by a
column pattern.

\begin{lemma}[Routing]\label{earlier:r:lem:routing}
Let \(U,V:[0,N)\times[0,C)\to\{0,1\}\). Switch values satisfying the
routing equations with \(D_0=U\) and \(D_{2k_N-1}=V\) exist if and only
if some permutation \(\pi\) of \([0,N)\) satisfies
\(V(\pi(x),c)=U(x,c)\) for all \(x\) and \(c\).
\end{lemma}
\begin{proof}
Soundness. The switch equations make each \(W_s(x,c)\) a bit, constant
in \(c\), and equal at the two lanes \(x\) and \(x\oplus2^{d_s}\). So the
map \(p_s\) sending \(x\) to \(x\) or to \(x\oplus2^{d_s}\), according to
the switch, is an involution, and the routing equation says
\(D_{s+1}(p_s(x),c)=D_s(x,c)\) for every column. Composing the
\(2k_N-1\) involutions gives \(\pi\); whole records move together, so
no record is lost or duplicated, even when values repeat.

Completeness, by induction on \(k_N\), following
Beneš~\cite{Ben65}. For \(k_N=1\)
one switch realizes either permutation of two lanes. For \(k_N\ge2\)
form a bipartite multigraph whose left vertices are the source pairs
\(\{x,x\oplus1\}\), whose right vertices are the destination pairs
\(\{y,y\oplus1\}\), and which has one edge from the pair of \(x\) to the
pair of \(\pi(x)\) for each \(x\). Every vertex has degree two, so the
multigraph is a union of even cycles, and its edges can be colored
\(0\) and \(1\) so that each vertex meets both colors. The first switch
of each source pair sends its record of color \(b\) to the lane whose
lowest bit is \(b\), and the last switch of each destination pair
receives its color-\(b\) record from that lane. For each color the
records induce a permutation of the remaining \(k_N-1\) lane bits,
which is routed by induction inside the two subnetworks, whose
dimension sequence is \(1,\ldots,k_N-1,\ldots,1\). Together with the
first and last stages this is the sequence
\(d_s=\min(s,2k_N-2-s)\). Copy every switch value to all columns and
propagate the records through the network; then all equations hold.
\end{proof}

\emph{Input subspace.} Let \(I\) have offset \(E(T,0,0,0,0)\) and the basis
\(e_{k_N+\ell_S},\ldots,e_{k_N+\ell_S+\omega-1}\),
\(e_0,\ldots,e_{m-\omega-1}\), in this order. Since \(c<N_A\le T\) and \(T\) is a single lane bit,
\(T+c=T\oplus c\), and the \(i\)-th element of \(I\) is the point of
bit \(i\bmod w\) of the value of \(\iota_{\floor{i/w}}\), because
\(\mathit{va}\) has field tag~\(0\). Thus the input clause of
\cref{earlier:r:def:acsp} writes \(\bar x\) directly into the input records. When
(B5) is compiled into \(P\), only the equations prescribing these
records' value bits are omitted; the input clause enforces them
instead. All other equations of (B5), including the zero values of the
dummy records, are retained.

\begin{theorem}[The affine instance]\label{earlier:r:thm:embedding}
Let \(\mathcal I(M,\tau,\bar x)\) be the affine constraint instance with
an irreducible polynomial of the degree \(f\) of \cref{earlier:r:lem:acsp-degree},
the subspace \(H\), the neighbors
above, the constraint polynomial \(P\) given by the sum, over the
equations of the event trace other than the input-value equations of
(B5), the routing equations and the bit checks, of the equation's
selector times the equation, and the input subspace
\(I\). It accepts \(\bar x\) if and only if
\(\mathrm{ET}(M,\tau,\bar x)\) is satisfiable.
\end{theorem}
\begin{proof}
At every point of \(H\) at most one selector is nonzero, because
equations have distinct host points or disjoint lane conditions. So
\(P(\alpha,A(\alpha+b_1),\ldots)\) vanishes on \(H\) exactly when every
equation holds for the table of values of \(A\) on \(H\); all points read
lie in \(H\), since every \(b_\nu\in H\). (\(\Rightarrow\)) Given \(A\),
read the slots of the trace from its values on \(H\). Every compiled
equation of the event trace holds, the bits are bits, and by \cref{earlier:r:lem:routing}
the network records at the last stage are a permutation of those at
stage~\(0\); hence so are their record fields, which is (P). The input
clause gives the values of the input records, the remaining equations
of (B5). (\(\Leftarrow\)) Given a
satisfying assignment of the event trace, place its values, set unused
network padding to \(0\), choose switches by \cref{earlier:r:lem:routing} for the
permutation of (P), set all other points to \(0\), and let \(A\) be the
interpolation polynomial of this table, of degree less than
\(\abs H\). Every equation and the input clause hold.
\end{proof}

\subsection{Degree, field, and generator}

\begin{lemma}[Degree and field]\label{earlier:r:lem:acsp-degree}
The constraint polynomial \(P\) of \cref{earlier:r:thm:embedding}, as emitted, has
weighted degree at most \(d_B=(m_H+4)2^{m_H}\). Let
\(f=m_H+\ceil{\lgx(4(m_H+4))}\). Then \(4d_B\le2^f\), \(m_H+2\le f\)
and \(2^f<8(m_H+4)2^{m_H}\). For the machine \(M_n\) of
\cref{earlier:r:lem:tm-ram} with \(T_n=2^{\tau_n}\),
\[
 2^f=O(T(n)\log^3T(n)),\qquad f=\log T(n)+O(\log\log T(n)),\qquad
 c_N=O(\log T(n)),
\]
with constants depending only on \(V\).
\end{lemma}
\begin{proof}
Every term is an emitted pattern, of degree at most \(m_H2^{m_H-1}\) in
\(X\), times an equation of degree at most two in the \(Y_\nu\); the
position constants of (B5) are single coordinates \(\lambda_a(X)\). A
term therefore has weight at most
\(m_H2^{m_H-1}+2(2^{m_H}-1)\le d_B\), and a sum of terms has no monomial
outside theirs. The bound holds for the polynomial as emitted, not only
on \(H\); this is why conditions are merged into single patterns before
emission. The choice of \(f\) gives \(2^f\ge4(m_H+4)2^{m_H}=4d_B\) and
\(2^f<8(m_H+4)2^{m_H}\). For \(M_n\), \(K\), \(L\) and \(F_{\rm tags}\)
depend only on \(V\), while \(w_n\) and \(k_N=\tau_n+1\) are
\(O(\log T(n))\) and \(T_n=O(T(n))\). Hence
\(\abs H=16N2^{\ell_S} C=O(T(n)\log^2T(n))\), and the bounds follow.
\end{proof}

\begin{lemma}[Generator]\label{earlier:r:lem:generator}
Fix \(V\). There is an algorithm that, given \(x\in\{0,1\}^n\) with
\(n\ge2\) and an integer \(t\ge n\) in binary, outputs \(\bar x\) and
the instance \(\mathcal I_x=\mathcal I(M_n,\tau_n,\bar x)\) of
\cref{earlier:r:thm:embedding}, built with \(t\) in place of \(T(n)\) in
\cref{earlier:r:lem:tm-ram}, in \(\poly(n,\log t)\) steps. The description of
\(\mathcal I_x\) has \(\poly(\log t)\) bits, and its constraint circuit
has \(O(\log^3t)\) gates, with constants depending only on \(V\).
\end{lemma}
\begin{proof}
The algorithm computes the numbers \(w_n\), \(\tau_n\), \(m\), \(k_N\),
\(\ell_S\), \(d\), \(q\), \(m_H\) and \(f\), all of \(O(\log t)\) bits, and
writes \(\bar x\). It obtains \(\mathrm{irr}\) from Shoup's deterministic
algorithm~\cite[Corollary~3.2]{Sho88}, which runs in \(\poly(f)\)
time~\cite[Algorithm~17]{BCGT13}, and solves the \(m_H\times m_H\) system
for the coefficients of the \(\lambda_a\) with \(O(m_H^3)\) field
operations. It emits the neighbor offsets, which are sums of basis
elements; the guard patterns, at most
\(4C(k_N+3)+4k_N(q+3)+4=O(\log^2t)\) of them, each with at most \(m_H\)
coordinates; and, for each equation family and bit index, the selector
times the equation, whose reads are the offsets of its catalog
entries. The code-dependent constants come from the fixed code of
\(M_n\) with \(n\), \(t\) and \(T_n-8\) inserted. Each pattern costs
\(O(m_H)\) gates and each equation \(O(L)\) gates, which gives
\(O((m_H+K+L)\cdot\log^2t+m_H^2+c_N)=O(\log^3t)\) gates. Every loop
runs over bit indices, code positions, registers, coordinates, neighbors
or patterns, never over the points of \(H\), the steps of the run, or a
witness; each field operation costs \(\poly(f)\) steps.
\end{proof}

\subsection{The short PCP}

\begin{theorem}[\Cref{imp:transfer-short-pcp} for multitape verifiers]
\label{earlier:r:thm:short-pcp-ours}
Let \(V(x,y)\) be a deterministic multitape machine and \(T:\N\to\N\) a
function with \(T(n)\ge n\), such that \(V\) halts within \(T(n)\) steps
on every \(x\in\{0,1\}^n\) and every \(y\), and \(T(n)\) is computable
from \(n\) in \(\poly(n,\log T(n))\) steps. Then, at every input length \(n\ge1\), the conclusion of
\cref{imp:transfer-short-pcp} holds for \(V\) and \(T\). In particular it holds
for \(\mathcal V_G\) and \(T_G\) of \cref{earlier:r:lem:succinct-verifier}.
\end{theorem}

\emph{Idea of proof.} \Cref{earlier:r:lem:tm-ram} turns \(V\) into a random-access
machine whose run has exactly \(T_n\) configurations, and
\cref{earlier:r:thm:trace,earlier:r:thm:embedding} turn the question whether that machine
accepts into an affine constraint instance (Step~1). The PCP verifier
of~\cite{BCGT13} for such
instances repeats its proximity sub-verifiers, which costs randomness;
we run each of them once (Step~2). The resulting verifier \(V_0\) is
sound only in a weak sense: it rejects every proof of a false instance
with probability at least \(1/\poly(\log T)\). A small-bias set amplifies this to soundness
\(1/n\) without new random bits (Step~3). Finally, the verifier is
written as a circuit whose proof reads are oracle gates (Step~4). The
governing bounds are
\[
 \begin{aligned}
 r_0&\le\log T+O(\log\log T),&\qquad
 \sigma_0&\ge\frac1{\poly(\log T)},\\
 \abs S&=\poly(n,\log T),&
 \ell&=\log T+O(\log\log T).
 \end{aligned}
\]
for the randomness, the base rejection probability, the size of the
small-bias set, and the number of inputs of \(C^O_x\). Every step runs
in time \(\poly(n,\log T)\).

\begin{proof}
For \(n=1\), use the zero-input constant-false circuit, as checked
below in ``The exact Lean interface''; the fixed resource constants
can cover this one length. Now fix \(n\ge2\) and
\(x\in\{0,1\}^n\). Write \(T=T(n)\), and take \(w_n\),
\(\bar x\), \(T_n\) and \(M_n\) from \cref{earlier:r:lem:tm-ram}. The construction
needs \(T\) explicitly, because \(w_n\), \(T_n\) and the code of \(M_n\)
depend on it. By hypothesis \(T\), and hence \(w_n\) and \(T_n\), are
computed in \(\poly(n,\log T)\) steps. All logarithms below are base two, and
\(\poly(n,\log T)\) bounds have polynomials depending only on \(V\) and
\(T\).

\emph{Step 1: reduction to an affine constraint instance.} Let
\(T_n=2^{\tau_n}\). By \cref{earlier:r:cor:trace-verifier,earlier:r:thm:embedding},
\(V(x,y)\) accepts for some \(y\) of length \(T\) if and only if the
affine constraint instance \(\mathcal I_x=\mathcal I(M_n,\tau_n,\bar x)\)
accepts \(\bar x\); here \(\abs{\bar x}=2^m\) with
\(\omega_n\le m\le\tau_n\) for \(w_n=2^{\omega_n}\), by
\cref{earlier:r:lem:tm-ram}. \Cref{earlier:r:lem:generator} computes \(\bar x\) and
\(\mathcal I_x\) from \(x\) and \(T\) in \(\poly(n,\log T)\) steps. Its
neighbors are translations, \(m_H+2\le f\), its input subspace has
dimension \(\lgx\abs{\bar x}\), with \(1\le\lgx\abs{\bar x}\le\tau_n<m_H\),
and, by \cref{earlier:r:lem:acsp-degree},
\[
 f=\log T+O(\log\log T),\qquad 1\le c_N=O(\log T).
\]

\emph{Step 2: the verifier with single sub-verifier runs.} Let \(V_0\)
be the verifier of~\cite[Construction~8.5]{BCGT13} for
\(\mathcal I_x\), with two changes. Each of the amplified
proximity verifiers \(V_{\rm aRS}\) and \(V_{\rm aVRS}\) in its Steps~2,
3 and~5 is replaced by a single run of the corresponding strong
proximity verifier \(V_{\rm RS}\) or \(V_{\rm VRS}\); and the test of
\(p_1\) uses the degree bound \(d_H\) below in place of the degree \(d\)
of~\cite[Construction~8.5]{BCGT13}. All steps read
prefixes of one random string, as in~\cite[Algorithm~9]{BCGT13}. The
proof format is that of~\cite[Definition~8.4]{BCGT13}: five tables of
elements of \(F\), namely an assignment table \(p_0\) and a composition
table \(p_1\), both indexed by \(F\), and three proximity proofs. We
encode them in binary with \(f\) bits per element. The test of \(p_1\)
uses the degree bound \(d_H=2^{f-2}=\abs F/4\), which \(V_0\) computes
from \(f\) alone; it bounds the degree \(d_B\) of \cref{earlier:r:lem:acsp-degree},
since \(4d_B\le2^f\), and it is less than \(2^f\). We use the following
properties.
\begin{enumerate}[label=(\alph*)]
\item \emph{Randomness and length.} The sub-verifiers use the
parameter tuple \((\eta,\kappa_0,\gamma_{\rm rep},\mu_{\rm rep})=(2,4,0,1)\)
fixed before \cref{earlier:r:lem:rs-recursion}. By~\cite[Lemmas~C.2--C.5]{BCGT13},
one run of \(V_{\rm RS}\) or \(V_{\rm VRS}\) on a domain of dimension
\(f\) uses at most \(\max\{\mathrm{rand}_=(f),f\}\) random bits, and its
proof has at most \(2^f+(2^\eta+1)\bigl(2^f+2\,\mathrm{len}_=(f)\bigr)\)
field elements. For the second bound, the degree bounds used by
\(V_0\) are less than \(2^f\), so the number of blocks
in~\cite[Lemma~C.3]{BCGT13} is at most \(2^\eta\). Since \(f\ge1\) and
\(\eta=2\), \cref{earlier:r:lem:rs-recursion} bounds these quantities by
\(f+3\log(2f)\) random bits and \(16\cdot2^ff^3\) field elements. The
consistency test uses \(f\) bits. As the steps share their random
string, \(V_0\) uses
\[
 r_0\le f+3\log(2f)=\log T+O(\log\log T)
\]
random bits; compare~\cite[Lemma~C.8]{BCGT13}, where each term is
multiplied by the number of repetitions. The proof consists of two
tables of \(2^f\) field elements and three proximity proofs. So the
binary proof has \(L_\pi\) bits, where
\[
 L_\pi\le f\bigl(2\cdot2^f+3\cdot16\cdot2^ff^3\bigr)\le64\cdot2^ff^4,
 \qquad \log L_\pi=\log T+O(\log\log T).
\]
\item \emph{Time and queries.} \(V_0\) runs in time
\[
 \poly(n,\log T).
\]
We bound the parts of its work in turn.
\begin{itemize}
\item \emph{Field operations.} The field \(F\) is represented by an
irreducible polynomial that Shoup's deterministic algorithm finds in
\(\poly(f)\) steps~\cite[Algorithm~17]{BCGT13}. So each field operation
takes \(\poly(f)\) steps.
\item \emph{Parameters.} The instance \(\mathcal I_x\) is generated in
\(\poly(n,\log T)\) steps (\cref{earlier:r:lem:generator}). Its constraint
circuit, which the consistency test evaluates once, has
\(O(\log^3T)\) gates, and its \(c_N\) neighbors are translations.
\item \emph{Sub-verifier runs.} A single run of \(V_{\rm RS}\) or
\(V_{\rm VRS}\) takes \(\poly(f)\) field operations. Its recursion
roughly halves the dimension of the current affine subspace at each
level and does \(\poly(f)\) field operations per
level~\cite[Algorithms~12--16]{BCGT13}.
\item \emph{The instance.} The only step whose cost depends on
\(\abs{\bar x}\) is the evaluation of the low-degree extension \(A_x\)
of the instance at the queried points (Step~5
of~\cite[Construction~8.5]{BCGT13}). By Lagrange interpolation over the
input subspace \(I\) of size \(\abs{\bar x}\), each
evaluation takes \(O(\abs{\bar x}^2)\) field operations, and
\(\abs{\bar x}=O(n+\log T)\).
\end{itemize}
The queries of \(V_0\) are computed from the random string alone, by
the non-adaptive decompositions of the single-run
verifiers~\cite[Lemmas~12.16--12.20]{BCGT13}, combined as
in~\cite[Lemma~12.23]{BCGT13}.
\item \emph{Completeness.} Suppose that \(\mathcal I_x\) accepts
\(\bar x\) with an assignment polynomial \(A\), and let \(A_x\) be the
polynomial of degree less than \(2^m\) that interpolates \(\bar x\) on
\(I\). Let \(\pi\) consist of the tables of \(A\) and of \(B_A\) on \(F\)
together with the proximity proofs of~\cite[Construction~8.3]{BCGT13},
with \(\pi_1\) computed for the degree bound \(d_H\).
Each table tested by \(V_0\) is a codeword of the code tested there:
\(A\) has degree less than \(2^{m_H}\), \(B_A\) has degree at most
\(d_B\le d_H\) and vanishes on \(H\), and \(A-A_x\) has degree less than \(2^{m_H}\) and
vanishes on \(I\). By the perfect completeness of the strong proximity
verifiers~\cite[Definition~4.15 and Theorem~11.1]{BCGT13}, every single
run accepts on every random string. The consistency test accepts, since
\(p_1(\alpha)=B_A(\alpha)=P(\alpha,A(N_1(\alpha)),\ldots)\). So \(V_0\)
accepts \(\pi\) with probability one.
\item \emph{Base soundness.} Suppose that \(\mathcal I_x\) does not
accept \(\bar x\), and let \(\pi\) be any proof. Put
\(\delta=1/(8c_N)\); here \(c_N\ge1\), because the identity is a
neighbor. We distinguish four cases. In Cases~1--3, one of the following
holds: \(p_0\) has relative distance at least \(\delta\) from every
polynomial of degree less than \(2^{m_H}\) (Case~1); \(p_1\) has relative
distance at least \(1/8\) from every polynomial of degree at most \(d_H\)
that vanishes on \(H\) (Case~2); or \(p_0-A_x\) has relative distance at
least \(\delta\) from every polynomial of degree less than \(2^{m_H}\)
that vanishes on \(I\) (Case~3). By strong soundness
\cite[Definition~4.15]{BCGT13} and~\cite[Theorem~11.1]{BCGT13}, a single
run of \(V_{\rm RS}\) or \(V_{\rm VRS}\) on a table at relative distance
at least \(\delta'\) from its code rejects with probability at least
\(\delta'/(2+K_f)\). Here \(K_f=16f^{\log c}\), with \(c\) the constant of~\cite[Theorem~11.3]{BCGT13} at the tuple fixed in~(a); this is a weakening of the bound printed in~\cite[Theorem~11.1]{BCGT13}, and the companion's source mapping explains the choice.

In Case~4 none of these holds, so there are polynomials \(A\) and
\(A'\) of degree less than \(2^{m_H}\), with \(A'\) vanishing on \(I\),
and \(B'\) of degree at most \(d_H\), vanishing on \(H\), such that
\(p_0\) and \(A\), and \(p_0-A_x\) and \(A'\), each differ on at most a
\(\delta\) fraction of \(F\), and \(p_1\) and \(B'\) on at most a
\(1/8\) fraction. The polynomials \(A-A_x\) and \(A'\) agree wherever
both of the first two relations hold, so on at least
\((1-2\delta)\abs F\ge\frac34\abs F>2^{m_H}\) points, because
\(m_H+2\le f\). Hence \(A-A_x=A'\), and \(A\) agrees with \(\bar x\) on
\(I\). Since \(\mathcal I_x\) does not accept \(\bar x\), \(B_A\) does not
vanish on all of \(H\), so \(B_A\ne B'\), and these polynomials of degree
at most \(d_H=\abs F/4\) agree on at most \(\abs F/4\) points. The
consistency test of Step~4 draws \(\alpha\in F\) uniformly and accepts
only if \(p_1(\alpha)=P(\alpha,p_0(N_1(\alpha)),\ldots,
p_0(N_{c_N}(\alpha)))\). Every \(N_\nu\) is a bijection of \(F\), so with
probability at least \(1-1/8-c_N\delta=3/4\) we have
\(p_1(\alpha)=B'(\alpha)\) and \(p_0(N_\nu(\alpha))=A(N_\nu(\alpha))\) for
every \(\nu\), and then the test accepts only if
\(B'(\alpha)=B_A(\alpha)\). So it rejects with probability at least
\(3/4-1/4=1/2\).

Since \(V_0\) rejects whenever one of its steps rejects, and
\(\delta/(2+K_f)\) is the smallest of the four bounds, it rejects \(\pi\)
with probability at least
\[
 \sigma_0=\frac1{8c_N(2+K_f)}\ge\frac1{\poly(\log T)}.
\]
\end{enumerate}

\emph{Step 3: amplification.} Choose fixed positive integers
\(K_1,b_1\) such that \(1/(K_1f^{b_1})\le\sigma_0\).
The consistency test uses \(f\ge1\) random bits, so \(r_0\ge1\).
Choose the least nonnegative integer \(m\) such that
\(2^{2m}\ge r_0^2nK_1f^{b_1}\), and use the powering construction
in the proof of \cref{earlier:r:lem:small-bias} over \(\mathbb F_{2^m}\).
Its bias is at most \(\varepsilon=r_0/2^m\), with
\(\varepsilon^2\le1/(nK_1f^{b_1})\le\sigma_0/n\).
Thus \cref{earlier:r:lem:patch-r} applies with \(s=1/n\). Since \(n\ge2\),
minimality of \(m\) gives
\[
 \abs S=2^{2m}<4r_0^2nK_1f^{b_1}=\poly(n,\log T).
\]
Choosing \(m\) uses only integer comparisons; constructing the field
and listing the multiset also take \(\poly(n,\log T)\) time.
The resulting verifier \(V_1\) uses \(r_0\) random bits. It accepts the
proof of~(c) with probability one when \(\mathcal I_x\) accepts
\(\bar x\), and it accepts every proof with
probability at most \(1/n\) otherwise. It runs in \(\poly(n,\log T)\)
time.

\emph{Step 4: the oracle circuit.} Let
\[
 \ell=\max\{r_0,\ceil{\log L_\pi}\},\qquad\text{so that}\qquad
 \ell=\log T+O(\log\log T).
\]
For \(R\in\{0,1\}^\ell\), consider the computation of \(V_1\) on
\(\bar x\) with the random string formed by the first \(r_0\) bits of
\(R\). It runs in \(\poly(n,\log T)\) steps, and its proof bits are read
at addresses below \(L_\pi\), written with \(\ell\) bits. Convert this
computation into a circuit in the standard way, and replace each read of
a proof bit by an oracle gate on the \(\ell\) address bits. This gives
an oracle circuit of size \(\poly(n,\log T)\) that outputs the verdict
of \(V_1\) on the proof \(\pi_O\) whose bit at address \(a<L_\pi\) is
\(O(a)\). Let \(C^O_x\) be its negation. The algorithm \(A\) computes
\(T\), \(\mathcal I_x\) (\cref{earlier:r:lem:generator}), \(S\) and \(C^O_x\) in
\(\poly(n,\log T)\) time.

\emph{The properties of \(C^O_x\).} If \(V(x,y)\) accepts for some \(y\)
of length \(T\), Step~1 and property~(c) give a proof \(\pi\) that
\(V_1\) accepts on every random string. The oracle \(O\) with
\(O(a)=\pi_a\) for \(a<L_\pi\) and \(O(a)=0\) otherwise makes \(C^O_x\)
unsatisfiable. If \(V(x,y)\) rejects for every \(y\) of length \(T\),
\(\mathcal I_x\) does not accept \(\bar x\) (Step~1). So for every
oracle \(O\), \(V_1\) accepts \(\pi_O\) on at most a \(1/n\) fraction of
the strings in \(\{0,1\}^{r_0}\). Hence \(C^O_x\) has at least
\[
 2^\ell(1-1/n)
\]
satisfying assignments. Each oracle gate has \(\ell\) inputs, and
\(C^O_x\) has \(\ell\) inputs and size \(\poly(n,\log T)\).

Finally, \(\mathcal V_G\) and \(T_G\) satisfy the hypotheses by
\cref{earlier:r:lem:succinct-verifier}(i)--(ii).
\end{proof}

\begin{corollary}[Clock supplied as input]\label{earlier:r:cor:short-pcp-clock}
Fix a deterministic multitape verifier \(V\). There is a uniform
algorithm that, given \(x\in\{0,1\}^n\), \(n\ge2\), and a binary integer
\(t\ge n\), constructs the oracle circuit of
\cref{imp:transfer-short-pcp} with \(T(n)\) replaced by \(t\), provided \(t\)
bounds the running time of \(V(x,y)\) for every witness \(y\).
Its running time and circuit size are \(\poly(n,\log t)\), and its
randomness and oracle-address lengths are
\(\log_2t+O(\log\log t)\). Completeness and soundness have the same
quantifiers as in \cref{imp:transfer-short-pcp}, with witnesses of length \(t\).
In particular, for any bound \(T(n)\ge n\), the compiler may be given
\((x,T(n))\); it need not compute the function \(T\).
\end{corollary}
\begin{proof}
In the proof of \cref{earlier:r:thm:short-pcp-ours}, computing \(T(n)\) serves
only to obtain the integer used to choose the word width, padded
clock, reduction parameters, and verifier circuit. Supply that integer
as \(t\) and perform the same construction, clocking the simulation
after \(t\) steps. All these choices are effective in
\(\poly(n,\log t)\) time, with constants depending on the fixed
verifier. Under the stated promise, the clock changes no computation
on \(x\). The completeness and soundness arguments of Steps~1--4
are pointwise in \(x\) and this integer, so apply unchanged. The
construction does not test the promise.
\end{proof}

\paragraph{The exact Lean interface.}
The companion states \cref{imp:transfer-short-pcp} as
\texttt{Ordinary\allowbreak{}Oracle\allowbreak{}Circuit\allowbreak{}PCP}, with an ordinary multitape
machine as the constructor. We complete it as follows. Its verifier starts with
\(\operatorname{fr}(x)\) and \(\operatorname{fr}(y)\) on its two input
tapes, where
\(\operatorname{fr}(b_1\cdots b_j)=1b_1\cdots1b_j0\), followed by blank
zeros. The predicate \texttt{RunsInTime} bounds this actual computation
by \(t=T(n)\) on every witness list; the correctness clauses quantify
over witnesses of exactly \(t\) bits. For \(n\ge2\), use the construction
of \cref{earlier:r:cor:short-pcp-clock} with these initial tapes. The machine
\(M_n\) of \cref{earlier:r:lem:tm-ram} writes \(\operatorname{fr}(x)\) and
\(\operatorname{fr}(y)\) onto the two simulated tapes, where \(y\)
consists of the first \(t\) witness bits. This writes the frame
terminators at their prescribed positions and leaves zeros beyond them;
subsequent writes by the simulated verifier are preserved. Loading takes
\(O(n+t+w_n)\) steps, so, since \(n\le t\), \(T_n=O(t+\log t)\), and
\(M_n\) accepts exactly for the required \(t\)-bit witnesses. The
remaining PCP construction is unchanged.

For \(n=1\), output a zero-input, one-node constant-false circuit,
with \(\ell=0\) and oracle arity zero. It is unsatisfiable for every
oracle, so completeness holds; the soundness lower bound is
\(2^0(1-1/1)=0\). Since \(T(1)\ge1\), both resource inequalities hold
after enlarging the fixed constants. For \(n=0\), use the same fixed
output: the two resource and two correctness clauses apply only for
\(n\ge1\). Thus one constructor covers every input length.

The constructor receives
\(\operatorname{fr}(x)\operatorname{fr}(\operatorname{bits}(t))\), of
length \(2n+2\abs{\operatorname{bits}(t)}+2\), and never computes the
function \(T\). For \(n\ge2\), its bounds give, for fixed integers
\(K,e\),
\[
 2^\ell\le Kt\,\lceil\log_2(t+2)\rceil^e,
 \qquad
 \abs{C_x^O}\le K\bigl(n+\operatorname{natBitLength}(t)+1\bigr)^e.
\]
Here \(\operatorname{natBitLength}(t)=\lfloor\log_2t\rfloor+1\) for
\(t\ge1\), and \(\lceil\log_2(t+2)\rceil\) is the companion's
\texttt{logScale}. Inputs, constants and oracle gates count as nodes;
each oracle gate has \(\ell\) arguments, and the circuit outputs true
precisely on rejection, as above. A circuit with \(s\) nodes has a word
in the companion's self-delimiting format of length
\(O((s(\ell+1)+1)\log(s+\ell+2))\): write its arity, node count, node
tags and addresses, and output address in order. Hence both computing
and writing this word take \(\poly(n,\log(t+2))\) time on an ordinary
multitape machine, which is the constructor the interface asks for.

\begin{remark}[Non-adaptivity]\label{earlier:r:rem:short-pcp-nonadaptive}
The circuit \(C^O_x\) of the proof is non-adaptive: the address bits of
every oracle gate are computed from the inputs \(R\) alone
(\cite[Lemmas~12.16--12.20]{BCGT13}, combined as
in~\cite[Lemma~12.23]{BCGT13}, and \cref{earlier:r:lem:patch-r}(i)); the oracle
answers enter only the decision.
\end{remark}

\subsection{Proof of \texorpdfstring{\cref{earlier:r:lem:tm-ram}}{the random-access simulation}}

The construction has two parts. A random-access machine
(\cref{earlier:r:def:ram}) simulates \(V\) step by step. An idle phase then makes
the number of configurations exactly \(T_n\), as acceptance in time
\(T_n\) requires.

\emph{Idea of proof.} The machine \(M_n\) copies the input and the
witness onto the simulated tapes of \(V\), which it keeps in its memory,
interleaved. It then simulates \(V\) step by step. Each block of its
code pads its own length to a multiple of \(8\) and adds that length to
a step counter. Once \(V\) accepts, a counted idle loop pads the run, so
that the \texttt{out} instruction comes at exactly step \(T_n-1\).

\begin{proof}[Proof of \cref{earlier:r:lem:tm-ram}]
\emph{Layout.} Let \(V\) have \(k_V\) tapes, counting its input and
witness tapes. Memory cell \(2k_Vj+i\) of \(M_n\) holds cell \(j\) of
tape \(i\) of \(V\); with this stride, a head at position \(0\) is
recognized by one comparison. The tape symbols are coded as small
integers, and the blank as zero. The \(K=9+2k_V\) registers hold a step
counter, the verdict, the state of \(V\), the head positions, the
scanned symbols, and scratch values.

\emph{Loading the tapes.} First \(M_n\) reads from tape~A the words
that contain the first \(n\) bits of \(\bar x\), which are \(x\), and
writes \(x\) bit by bit onto the input tape of \(V\). It then reads from
tape~B the words that contain the first \(T(n)\) witness bits, and writes
these bits, \(y\), onto the witness tape of \(V\). (For the companion's
verifiers the two tapes receive the framed words
\(\operatorname{fr}(x)\) and \(\operatorname{fr}(y)\); see ``The exact
Lean interface'' above.) The integers \(n\) and \(T(n)\) are written
into the code. If a read returns \(\#\), \(M_n\) enters a one-instruction
loop that never reaches \texttt{out}. Loading takes
\(O(n+T(n)+w_n)\) steps.

\emph{Simulating \(V\).} One step of \(V\) is then simulated by a block
that loads the scanned symbols, forms the index of the pair (state,
scanned symbols), and branches through a chain of comparisons to one
leaf per pair. The leaf stores the written symbols, moves the heads and
sets the new state. A halted rejecting state, or a pair without a
transition, leads to the loop that never halts; the accepting state
records the verdict and starts the idle phase. A step costs \(O(1)\)
instructions for fixed \(V\), so the simulation ends within
\(O(n+w_n+T(n))=O(T(n)+w_n)\) steps.

\emph{Counting.} The code of \(M_n\) consists of blocks: the start, the
two loading loops, the simulation of one step of \(V\) with its leaves,
and the recording of the verdict. Every path through a block is padded
with \texttt{nop} instructions so that the number of steps it takes is
a multiple of~\(8\). On that path the block adds this number to a
counter register, which is zero in the initial configuration. Hence,
whenever a block ends, the counter equals the number of steps made so
far, and this number is a multiple of~\(8\). This at most multiplies
the running time by a constant.

\emph{Idling.} Once the verdict is recorded in a register, \(M_n\)
idles. A fixed code cannot hold a number of \texttt{nop}s that depends
on the run, so the padding is a counted loop: an idle block of exactly
\(8\) instructions. In order, they are:
\begin{itemize}
\item \texttt{add} \(8\) to the counter;
\item \texttt{cmp} the counter with the constant \(T_n-8\);
\item five \texttt{nop}s;
\item \texttt{jne} back to the start of the idle block if the comparison
found the counter different from \(T_n-8\); as the counter never exceeds
\(T_n-8\), this means smaller.
\end{itemize}
The counter takes only multiples of \(8\). Since \(T_n\) is a power of
two and at least \(16\), \(T_n-8\) is also a multiple of \(8\). So if the
counter is at most \(T_n-16\) on entry, the loop ends after the pass at
whose end exactly \(T_n-8\) steps have been made.

\emph{The verdict.} Then \(M_n\) executes six \texttt{nop}s and one
\texttt{je} on the verdict register. The \texttt{je} jumps to an
instruction \((\texttt{out},\textit{reject})\) if the verdict is
rejection, and otherwise falls through to an instruction
\((\texttt{out},\textit{accept})\). The instruction \texttt{out} takes a
constant argument, so this branch is needed. After these seven steps
exactly \(T_n-1\) steps have been made, and the \texttt{out} instruction
with the verdict is executed at step \(T_n-1\).

\emph{Constants.} A suitable constant \(c_V\ge16\) ensures three things.
\begin{itemize}
\item The simulation and the recording of the verdict end within
\(T_n-16\) steps.
\item \(\abs{\bar x}\le2(n+1)+2w_n\le T_n\).
\item \(2^{w_n}\ge2^{c_V}T(n)\) is larger than \(T_n\) and every
address.
\end{itemize}
The code of \(M_n\) is a fixed code with the integers \(n\), \(T(n)\)
and \(T_n-8\) written into it; its length \(L\) and the number \(K\) of
registers depend only on \(V\).

\emph{Runs.} If \(V(x,y)\) accepts for some \(y\in\{0,1\}^{T(n)}\), let
tape~B hold \(y\) padded with zeros to a multiple of \(w_n\), of length
at most \(T(n)+w_n\le T_n\). The computation of \(M_n\) on \(\bar x\) and
this tape, in which every read returns the next word, is an accepting
run of length \(T_n\). Conversely, no read of an accepting run returns
\(\#\), since that leads to the loop that never halts. So the run reads
the first words of \(\bar x\) and of a witness string, simulates \(V\) on
\(x\) and the first \(T(n)\) bits \(y\) of that string, and reaches the
accepting state; hence \(V(x,y)\) accepts.
\end{proof}


\subsection{Amplification and uniform proximity interfaces}\label{earlier:release-interfaces}
These are the arguments of \cite[Appendix~B]{NagCG}; they are
reproduced with notation adapted. Their published premises are the
STV reconstruction theorem and the two Chen--Williams results in
\cref{par:five-inputs}.

\paragraph{Full-range amplification interface.}
We justify the formulation of \cite[Lemma~3.9]{CLW20} used by the
companion, for every time-constructible \(S:\mathbb N\to\mathbb N\).
Here circuit size counts all DAG nodes, including input and constant
nodes. Following the construction of \cite[Lemmas~25 and~28]{STV01}
and the circuit argument of its Theorem~24, the polynomial
reconstruction theorem \cite[Theorem~29]{STV01} gives absolute constants
\(A\ge1\) and an integer \(a\ge1\) with the following consequence. An
\(s'\)-node circuit agreeing with the encoded function on more than
a \(1/2+\delta\) fraction yields an exact circuit for the original
\(n\)-input function with at most
\(A(n/\delta)^a s'+n+2\) nodes. The polynomial reconstruction
overhead includes fixing the decoder's randomness and conversion
to our fan-in-two basis; the last term allows input and constant
nodes. Choose an integer \(c>2a+2\), once and before \(S\).

Every \(n\)-input function has a circuit with fewer than
\(U(n)=2^{n+3}\) nodes. Indeed Shannon expansion combines two
\((n-1)\)-input circuits as
\((x\wedge f_1)\vee(\neg x\wedge f_0)\), adding at most five nodes.
Thus \(B(0)=1\) and \(B(n)\le2B(n-1)+5\) give
\(B(n)\le6\cdot2^n-5<U(n)\).
For all sufficiently large \(n\), construct the output as follows.
If \(S(n)=0\), use a constant function: the target size is zero and
no circuit has zero nodes. If \(1\le S(n)\le n^c\), use parity on
\(n+1\) bits. A circuit with at most \(S(n)^{1/c}\le n\) nodes
omits an input, so pairing assignments that differ only in that
input gives agreement exactly \(1/2<1/2+S(n)^{-1/c}\).
If \(S(n)\ge U(n)\), use a constant function, since no input
function satisfies the hardness premise.

In the remaining regime \(n^c<S=S(n)<U(n)\), put
\(\epsilon=S^{-1/c}\). Compute the least integer \(k\) such that
\(2^{ck}\ge S\), and use the dyadic accuracy
\(\delta=2^{-k-2}\); then
\(\epsilon/8<\delta\le\epsilon/4\).
All parameter choices use integer arithmetic. Apply the binary
encoder of \cite[Lemma~28]{STV01} to the truth table.
A circuit of size at most \(S^{1/c}\) with agreement at least
\(1/2+\epsilon\) would have agreement strictly greater than
\(1/2+\delta\). After increasing the absolute constant \(A\) to
absorb \(8^a\), reconstruction gives size at most
\[
 A n^a S^{(a+1)/c}+n+2<S.
\]
The strict inequality holds uniformly in this regime from a fixed
onset: dividing by \(S>n^c\) bounds the left side by
\[
 \frac{A}{n^{\,c-2a-1}}+\frac{n+2}{n^c}<1.
\]
This contradicts the premise and proves the required strict
agreement bound.

On the truth-table layout: in the construction of
\cite[Lemmas~25 and~28]{STV01} choose a binary extension field,
rounding its size up to the required scale, and round the number
of polynomial variables up. The evaluation domain is then
\(\mathbb F^m\); concatenation with a Hadamard block indexes the
binary codeword by a product of binary cubes. Its length is a
power of two, with no padding of the codeword, and remains
\(\poly(2^n,1/\delta)\). These roundings preserve the polynomial
reconstruction bound above. Since \(S<U(n)\), we have
\(1/\delta=2^{O(n)}\), so the output has \(O(n)\) input bits and
its entire truth table is computable in \(2^{O(n)}\) time.

Choose one onset \(n_0\ge2\) for these estimates, also large enough
that \(n^c<U(n)\) for \(n\ge n_0\). At each of the finitely many
lengths \(n<n_0\) with \(S(n)>0\), output parity on \(S(n)+1\)
bits instead. Since \(S(n)^{1/c}\le S(n)\), the same missing-input
argument gives hardness at every such length, for every supplied
function. The case \(S(n)=0\) is as above. These finitely many
choices can be hardcoded after \(S\) is fixed; the arity and
runtime bounds in the companion have onsets, while its hardness
clause holds at every length.

Finally the eventual runtime must hold even on easy input
functions and for very large \(S(n)\). Fix a time constructor
computing \(S(n)\) from unary \(n\) within \(C(S(n)+1)\) steps.
For \(n\ge n_0\), run it for only \(C(U(n)+1)\) simulated steps.
A timeout implies \(S(n)>U(n)\); a returned value at least \(U(n)\)
also selects the constant output. Otherwise the returned value
determines the cases above. Initialization, the capped simulation,
integer operations, and reading and writing the framed arities
and truth tables all take \(2^{O(n)}\) time. The usual
polynomial-overhead simulation into \texttt{TM2}, the multi-stack
machine model of the Lean mathematical library in which the companion
states this input, preserves this bound.
After increasing its exponent and onset, the bound is
\(2^{en}\) as required by the companion. The machine never tests
whether the input function is hard. Its constants and onsets may
depend on \(S\); the exponent \(c\) in the hardness conclusion
does not.

\paragraph{Uniform three-query PCPP and code interfaces.}
We spell out the uniformity and boundary conventions of the
companion's two inputs underlying \cref{imp:transfer-pcpp}.
For a fixed real \(0<\delta<1\), choose a positive rational
\(\eta<\delta\) and apply \cite[Lemma~22]{CW19} at proximity
\(\eta\) and soundness \(1/2\), followed by the constant-size clause
conversion of \cite[Appendix~A]{CW19} that proves
\cite[Lemma~24]{CW19}. Distance at least \(\delta\) from every
accepting word implies the source's strict distance condition at
\(\eta\). In that clause conversion,
round the common clause count up to a power of two and pad each
block with tautological three-literal clauses; use a dummy proof
coordinate if necessary and repeat literals in shorter clauses.
Uniform binary coins then select exactly three literals, and
padding decreases the rejection probability by at most a fixed
factor. The verifier is nonadaptive
(\cite[Definition~20]{CW19}):
for a fixed circuit \(C\) and coin string, its three queried
addresses and their signs are determined without oracle answers.
Enumerate all \(2^r=\poly(|C|)\) coin strings, list the proof
addresses they use, and sort and relabel the distinct addresses.
There are at most \(3\cdot2^r\) of them. Restricting an original
proof to these addresses preserves every verifier outcome;
conversely every trimmed proof extends arbitrarily at unused
addresses and has the same outcomes. Thus trimming preserves
both perfect completeness and soundness.

The verifier's uniform algorithms compute each address in
polynomial time, so the addresses have polynomial bit length.
Enumeration, comparison, relabeling, and emission of the
three-literal clauses are polynomial-time operations. Apply the
same restriction to the honest proof: the proof produced
by~\cite[Lemma~22]{CW19}, extended to the auxiliary bits of the
clause conversion by evaluating its constant-size circuits
(\cite[Appendix~A]{CW19}). Its uniform polynomial construction is
preserved. If \(m\) is the companion's circuit
node count, then \(n\le m\), and a topological node list has
ordinary description length \(O((m+1)\log(m+2))\).
The companion's balanced binary encodings also have polynomial
length in \(m+1\): each pairing multiplies bit length by a
constant and the list tree has logarithmic depth. Encoding and
decoding these descriptions take polynomial time. Consequently
the proof length and number of coin strings are bounded by
\(K(m+1)^e\), and the constructor and honest-proof algorithm
have polynomial-time TM2 realizations on the stated encodings.

For the code, the cited statement is \cite[Lemma~26]{CW19}: a
constant-rate binary linear code with positive minimum relative
distance, an efficient encoder, and an efficient decoder that corrects
a constant fraction of errors. The main result uses only the encoder
facts: linearity, so that each output coordinate is a parity of a
fixed support of message coordinates; constant rate; positive minimum
relative distance; and a polynomial-time encoder. These are what the
checked circuit and the distance argument after \cref{imp:transfer-pcpp}
use. The finitely many short lengths can use identity codes, with the
common distance constant decreased if needed. At \(n=0,1\) choose
arbitrary fixed interfaces, since the companion's code clauses
concern \(n\ge2\). The encoder has a polynomial-time TM2 realization.

The companion's code input is exactly these encoder facts, for
\(n\ge2\), derived from~\cite[Lemma~26]{CW19} as just described. The
companion also proves the stronger full code statement
of~\cite[Lemma~27]{CW19}, which is not needed here. Its three-block
construction
\(\operatorname{Enc}(x)=\operatorname{Enc}_0(x_1)\circ
\operatorname{Enc}_0(x_2)\circ\operatorname{Enc}_0(x_3)\), on the
thirds \(x_1,x_2,x_3\) of the message and with the code
\(\operatorname{Enc}_0\) of~\cite[Lemma~26]{CW19}, bounds every support
by \(\ceil{n/3}\le\floor{n/2}\) for \(n\ge2\); the stronger statement
also supplies a decoder. For it, the decoding constant is decreased together with the
distance at short lengths, and \(0<\rho<d<1\) are chosen below,
respectively, the guaranteed decoding radius and minimum relative
distance. Choosing \(\rho\) strictly inside the decoding guarantee
includes the closed ball of radius \(\rho\). Make the decoder total,
returning a fixed \(n\)-bit word on failure; this does not change its
value on a promised input. The total decoder also has a
polynomial-time TM2 realization. The support bound and the decoder are
retained as proved results; neither is required by the main encoder
interface or by the proof of \cref{imp:transfer-pcpp}.

